\documentclass[8pt]{beamer}
\usetheme{Madrid}
\usecolortheme{default}
\usepackage{amsmath, amssymb, graphicx, array, multirow, booktabs, xcolor, tikz}
\usetikzlibrary{shapes,arrows,positioning,fit,backgrounds}

% Compact font settings
\setbeamerfont{title}{size=\Large}
\setbeamerfont{subtitle}{size=\small}
\setbeamerfont{author}{size=\scriptsize}
\setbeamerfont{date}{size=\scriptsize}
\setbeamerfont{frametitle}{size=\large}
\setbeamerfont{framesubtitle}{size=\normalsize}
\setbeamerfont{enumerate item}{size=\footnotesize}
\setbeamerfont{itemize item}{size=\footnotesize}
\setbeamerfont{block title}{size=\footnotesize}

% Compact spacing
\setlength{\itemsep}{0.08ex}
\setlength{\parskip}{0.08ex}
\setlength{\leftmargini}{0.3cm}
\setbeamersize{text margin left=3mm, text margin right=3mm}
\addtobeamertemplate{frame begin}{\vspace{-0.4cm}}{}

% Colors
\definecolor{myblue}{RGB}{0,0,255}
\definecolor{myred}{RGB}{255,0,0}
\definecolor{mygreen}{RGB}{0,128,0}
\definecolor{myorange}{RGB}{255,165,0}
\definecolor{mypurple}{RGB}{128,0,128}
\definecolor{mygray}{RGB}{128,128,128}
\definecolor{mygold}{RGB}{255,215,0}

\title{Chapter 8: Supervised Learning - Model Performance Evaluation}
\subtitle{Comprehensive Tutorial Notes: Continuous Metrics, Classification Metrics, ROC/AUC, Threshold Selection}
\author{GARP RAI Exam Preparation}
\date{}

\begin{document}
	
	\begin{frame}
		\titlepage
	\end{frame}
	
	% ============================================================
	% SECTION 1: INTRODUCTION TO MODEL EVALUATION
	% ============================================================
	
	\begin{frame}{T1: Model Evaluation Purpose}
		\fontsize{6.5}{7.5}\selectfont
		\textbf{TOPIC: Primary Purpose of Model Performance Evaluation}
		
		\vspace{0.1cm}
		\textbf{DEFINITION:}
		\begin{itemize}
			\item Model performance evaluation is the systematic process of assessing how well a supervised learning model generalizes to new, unseen data
			\item It provides the technical basis to understand model accuracy and the right questions to ask when choosing appropriate performance metrics
		\end{itemize}
		
		\vspace{0.1cm}
		\textbf{KEY CONCEPTS:}
		\begin{enumerate}
			\item \textcolor{myblue}{Generalization Assessment:} How well does the model perform on data it has not seen during training?
			\item \textcolor{myblue}{Model Comparison:} Which model is best for the specific problem at hand?
			\item \textcolor{myblue}{Error Quantification:} How wrong can we afford to be, and what are the consequences when models are wrong?
			\item \textcolor{myblue}{Risk Management:} What are the financial and reputational implications of model errors?
		\end{enumerate}
		
		\vspace{0.1cm}
		\textbf{WHY IT MATTERS:}
		\begin{itemize}
			\item Models drive critical decisions in credit risk, fraud detection, AML compliance
			\item Poor models have significant financial and reputational consequences
			\item Regulatory requirements demand rigorous model validation
			\item Stakeholder confidence depends on reliable model performance
		\end{itemize}
		
		\vspace{0.1cm}
		\textbf{PRACTICAL EXAMPLE:}
		\begin{itemize}
			\item Greek bank estimated recovery rate at 65\%
			\item Actual recovery was 40\%
			\item Result: €120 million capital add-on, stock price slide, regulatory scrutiny
		\end{itemize}
		
		\textcolor{myred}{\textbf{EXAM TRAP:}} Model evaluation is about \textcolor{myblue}{generalization performance}, not training performance!
	\end{frame}
	
	\begin{frame}{T2: Continuous vs Classification Outputs}
		\fontsize{6.5}{7.5}\selectfont
		\textbf{TOPIC: Distinguishing Between Continuous and Discrete Output Metrics}
		
		\vspace{0.1cm}
		\textbf{DEFINITION:}
		\begin{itemize}
			\item Supervised learning models produce either continuous numerical outputs (regression) or discrete categorical outputs (classification)
			\item Each type requires fundamentally different evaluation approaches and metrics
		\end{itemize}
		
		\vspace{0.1cm}
		\textbf{CONTINUOUS OUTPUTS (REGRESSION):}
		\begin{itemize}
			\item \textbf{Predict:} Numbers such as VaR, market prices, net revenue, expected credit losses
			\item \textbf{Metrics:} MSE, RMSE, MAE, MAPE
			\item \textbf{Key Question:} How numerically close are predictions to actual values?
		\end{itemize}
		
		\vspace{0.1cm}
		\textbf{DISCRETE OUTPUTS (CLASSIFICATION):}
		\begin{itemize}
			\item \textbf{Predict:} Categories such as fraud vs legitimate, default vs paid off, approve vs decline
			\item \textbf{Metrics:} Accuracy, Precision, Recall, F1 Score, AUC
			\item \textbf{Key Question:} How well does the model distinguish between classes?
		\end{itemize}
		
		\vspace{0.1cm}
		\textbf{WHY THE DISTINCTION MATTERS:}
		\begin{itemize}
			\item Continuous metrics measure numerical accuracy (how far off is the prediction?)
			\item Classification metrics measure classification accuracy (how many correct classifications?)
			\item Different types of errors have different costs in each context
		\end{itemize}
		
		\textcolor{myred}{\textbf{EXAM TRAP:}} \textcolor{myblue}{Continuous vs discrete} outputs require fundamentally different evaluation approaches!
	\end{frame}
	
	\begin{frame}{T3: Consequences of Poor Modeling}
		\fontsize{6.5}{7.5}\selectfont
		\textbf{TOPIC: Real-World Consequences of Poor Model Performance}
		
		\vspace{0.1cm}
		\textbf{THE CHAIN OF CONSEQUENCES:}
		\begin{enumerate}
			\item \textcolor{myblue}{Financial Impact:}
			\begin{itemize}
				\item Capital add-ons (e.g., €120 million for Greek bank)
				\item Shareholder dilution
				\item Underpricing of risk in new loan originations
				\item Future losses in portfolio
			\end{itemize}
			
			\item \textcolor{myblue}{Operational Impact:}
			\begin{itemize}
				\item Deep-dive reviews of entire credit risk modeling suite
				\item Significant internal resource consumption
				\item Potential systemic weaknesses requiring costly overhauls
			\end{itemize}
			
			\item \textcolor{myblue}{Reputational Impact:}
			\begin{itemize}
				\item Lingering drag on investor confidence
				\item Damaged relationships with regulators
				\item Market confidence erosion
			\end{itemize}
			
			\item \textcolor{myblue}{Regulatory Impact:}
			\begin{itemize}
				\item ECB supervisory review actions
				\item Increased scrutiny of all models
				\item Potential fines and sanctions
			\end{itemize}
		\end{enumerate}
		
		\textcolor{myred}{\textbf{EXAM TRAP:}} Poor modeling has \textcolor{myblue}{real financial, operational, and reputational consequences}!
	\end{frame}
	
	\begin{frame}{T4: Test Sample Importance}
		\fontsize{6.5}{7.5}\selectfont
		\textbf{TOPIC: Why Compute Metrics on Test Sample}
		
		\vspace{0.1cm}
		\textbf{DEFINITION:}
		\begin{itemize}
			\item The test sample comprises observations not used in determining parameter estimates or tuning hyperparameters
			\item It provides an unbiased estimate of how the model will perform on new, unseen data
		\end{itemize}
		
		\vspace{0.1cm}
		\textbf{THE OVERFITTING PROBLEM:}
		\begin{itemize}
			\item Models learn training data patterns, including noise and random fluctuations
			\item Performance on training data is optimistically biased
			\item Training performance is not a true test of generalization
			\item A model that memorizes training data may fail completely on new data
		\end{itemize}
		
		\vspace{0.1cm}
		\textbf{DATA SPLIT PURPOSES:}
		\begin{itemize}
			\item \textcolor{myblue}{Training Sample:} Estimate model parameters
			\item \textcolor{myblue}{Validation Sample:} Tune hyperparameters and select model architecture
			\item \textcolor{myblue}{Test Sample:} Evaluate final model performance (unseen data)
		\end{itemize}
		
		\vspace{0.1cm}
		\textbf{KEY INSIGHT:}
		\begin{itemize}
			\item We are generally more interested in validation and test performance
			\item Examining accuracy on training data is not a true test of performance
			\item Test sample performance is what matters in practice
		\end{itemize}
		
		\textcolor{myred}{\textbf{EXAM TRAP:}} Always evaluate on \textcolor{myblue}{test sample} for true generalization performance!
	\end{frame}
	
	\begin{frame}{T5: Training vs Test Gap}
		\fontsize{6.5}{7.5}\selectfont
		\textbf{TOPIC: Interpreting Training-Test Performance Gap}
		
		\vspace{0.1cm}
		\textbf{DEFINITION:}
		\begin{itemize}
			\item The difference between model performance on training data versus test data
			\item A key diagnostic for detecting overfitting
		\end{itemize}
		
		\vspace{0.1cm}
		\textbf{PERFORMANCE PATTERNS:}
		\begin{itemize}
			\item \textcolor{mygreen}{No Gap:} Training performance $\approx$ Test performance (good generalization)
			\item \textcolor{myorange}{Small Gap:} Slight overfitting (acceptable in practice)
			\item \textcolor{myred}{Large Gap:} Significant overfitting (problematic, needs remediation)
		\end{itemize}
		
		\vspace{0.1cm}
		\textbf{WHY THE GAP OCCURS:}
		\begin{itemize}
			\item Model memorizes training patterns rather than learning generalizable relationships
			\item Learns noise instead of signal
			\item Too complex for the amount of training data available
			\item Insufficient regularization
		\end{itemize}
		
		\vspace{0.1cm}
		\textbf{EXAMPLE (LendingClub):}
		\begin{itemize}
			\item Logistic Regression: Training accuracy = 84.2\%, Test accuracy = 71.3\%
			\item Neural Network: Training accuracy = 84.2\%, Test accuracy = 70.1\%
			\item Both show significant gap indicating overfitting
		\end{itemize}
		
		\vspace{0.1cm}
		\textbf{REMEDIATION:}
		\begin{itemize}
			\item Remove irrelevant features
			\item Apply regularization (L1, L2, dropout)
			\item Get more training data
			\item Use cross-validation
		\end{itemize}
		
		\textcolor{myred}{\textbf{EXAM TRAP:}} Large train-test gap = \textcolor{myblue}{overfitting}!
	\end{frame}
	
	% ============================================================
	% SECTION 2: CONTINUOUS OUTPUT METRICS
	% ============================================================
	
	\begin{frame}{T6: Mean Squared Error (MSE)}
		\fontsize{6.5}{7.5}\selectfont
		\textbf{TOPIC: Mean Squared Error - Definition and Calculation}
		
		\vspace{0.1cm}
		\textbf{FORMULA:}
		\[
		MSE = \frac{1}{N}\sum_{i=1}^N (y_i - \hat{y}_i)^2
		\]
		where:
		\begin{itemize}
			\item $y_i$ = true value of outcome for observation $i$
			\item $\hat{y}_i$ = model prediction for observation $i$
			\item $N$ = number of observations in test sample
		\end{itemize}
		
		\vspace{0.1cm}
		\textbf{WHAT IT MEASURES:}
		\begin{itemize}
			\item Average of all squared forecast errors
			\item Squaring removes sign (under- and over-estimation penalized symmetrically)
			\item \textcolor{myblue}{Heavily penalizes large errors} (tail risk focus)
		\end{itemize}
		
		\vspace{0.1cm}
		\textbf{KEY CHARACTERISTICS:}
		\begin{itemize}
			\item Always non-negative
			\item Lower is better (most accurate model has lowest MSE)
			\item Units are squared (e.g., dollars²) - difficult to interpret intuitively
			\item Sensitive to outliers
		\end{itemize}
		
		\vspace{0.1cm}
		\textbf{EXAMPLE (Salary Prediction):}
		\begin{itemize}
			\item Model: $\hat{y} = 16.88 + 1.06 \times \text{experience}$
			\item Sum of squared errors = 431.02
			\item N = 10
			\item MSE = 431.02 / 10 = 43.1
		\end{itemize}
		
		\textcolor{myred}{\textbf{EXAM TRAP:}} MSE uses \textcolor{myblue}{squared errors} - heavily penalizes large deviations!
	\end{frame}
	
	\begin{frame}{T7: MSE Managerial Question}
		\fontsize{6.5}{7.5}\selectfont
		\textbf{TOPIC: What Managerial Question Does MSE Answer?}
		
		\vspace{0.1cm}
		\textbf{MANAGERIAL QUESTION:}
		\begin{itemize}
			\item "How do we penalize big errors in our model?"
		\end{itemize}
		
		\vspace{0.1cm}
		\textbf{MANAGERIAL FOCUS:}
		\begin{itemize}
			\item Highlights tail risk
			\item Emphasis on squared errors heavily penalizes large deviations
			\item Sensitive to costly outliers
			\item Important for risk management
		\end{itemize}
		
		\vspace{0.1cm}
		\textbf{WHY THIS MATTERS:}
		\begin{itemize}
			\item Large errors have significant consequences in risk management
			\item Risk managers care about tail events (extreme losses)
			\item Regulatory focus on extreme outcomes and capital adequacy
			\item MSE aligns with regulatory capital frameworks
		\end{itemize}
		
		\vspace{0.1cm}
		\textbf{REGULATORY PERSPECTIVE:}
		\begin{itemize}
			\item ECB's Guide to Internal Models recommends metrics capturing tail risk
			\item Ensures large, infrequent losses are not overlooked
			\item RMSE reported alongside MAE for complete picture
		\end{itemize}
		
		\textcolor{myred}{\textbf{EXAM TRAP:}} MSE is for \textcolor{myblue}{penalizing big errors} - risk managers love it!
	\end{frame}
	
	\begin{frame}{T8: Root Mean Square Error (RMSE)}
		\fontsize{6.5}{7.5}\selectfont
		\textbf{TOPIC: Root Mean Square Error - Definition and Interpretation}
		
		\vspace{0.1cm}
		\textbf{FORMULA:}
		\[
		RMSE = \sqrt{MSE} = \sqrt{\frac{1}{N}\sum_{i=1}^N (y_i - \hat{y}_i)^2}
		\]
		
		\vspace{0.1cm}
		\textbf{KEY CHARACTERISTICS:}
		\begin{itemize}
			\item Same units as target variable (e.g., dollars)
			\item More interpretable than MSE
			\item Gives sense of average error magnitude
			\item Still sensitive to outliers due to squaring before square root
		\end{itemize}
		
		\vspace{0.1cm}
		\textbf{COMPARISON TO MSE:}
		\begin{itemize}
			\item MSE: Units are squared (e.g., dollars²) - hard to interpret
			\item RMSE: Units are original (e.g., dollars) - easy to interpret
			\item Both penalize large errors disproportionately
		\end{itemize}
		
		\vspace{0.1cm}
		\textbf{EXAMPLE (Salary Prediction):}
		\begin{itemize}
			\item MSE = 43.1
			\item RMSE = $\sqrt{43.1} = 6.57$
			\item Interpretation: "On average, predictions are off by about 6.57 thousand dollars"
		\end{itemize}
		
		\vspace{0.1cm}
		\textbf{CAVEAT FOR MANAGERS:}
		\begin{itemize}
			\item Unusually large loss on single facility could disproportionately inflate RMSE
			\item Need to ask: "Is high RMSE driven by systemic issue or one or two extreme events?"
		\end{itemize}
		
		\textcolor{myred}{\textbf{EXAM TRAP:}} RMSE is \textcolor{myblue}{easier to interpret} than MSE because it's in same units!
	\end{frame}
	
	\begin{frame}{T9: RMSE Managerial Question}
		\fontsize{6.5}{7.5}\selectfont
		\textbf{TOPIC: What Managerial Question Does RMSE Answer?}
		
		\vspace{0.1cm}
		\textbf{MANAGERIAL QUESTION:}
		\begin{itemize}
			\item "What does the typical error feel like in our primary unit of currency (e.g., dollars)?"
		\end{itemize}
		
		\vspace{0.1cm}
		\textbf{MANAGERIAL FOCUS:}
		\begin{itemize}
			\item Same units as target variable (e.g., dollars)
			\item Directly interpretable for non-technical stakeholders
			\item "On average, our predictions are off by RMSE dollars"
			\item Helps CEOs and board members understand model performance
		\end{itemize}
		
		\vspace{0.1cm}
		\textbf{WHY THIS MATTERS:}
		\begin{itemize}
			\item Connects statistical measure to business impact
			\item Board members can relate to dollar figures
			\item Enables comparison across different models
			\item Facilitates communication with non-technical audiences
		\end{itemize}
		
		\vspace{0.1cm}
		\textbf{CAVEAT:}
		\begin{itemize}
			\item Still sensitive to outliers (squaring before square root)
			\item An unusually large loss on single facility could disproportionately inflate RMSE
			\item May mask generally good performance across rest of portfolio
			\item Need to investigate if high RMSE is systemic or from extreme, explainable events
		\end{itemize}
		
		\textcolor{myred}{\textbf{EXAM TRAP:}} RMSE answers \textcolor{myblue}{"what does typical error feel like?"}
	\end{frame}
	
	\begin{frame}{T10: Mean Absolute Error (MAE)}
		\fontsize{6.5}{7.5}\selectfont
		\textbf{TOPIC: Mean Absolute Error - Definition and Advantages}
		
		\vspace{0.1cm}
		\textbf{FORMULA:}
		\[
		MAE = \frac{1}{N}\sum_{i=1}^N |y_i - \hat{y}_i|
		\]
		
		\vspace{0.1cm}
		\textbf{KEY CHARACTERISTICS:}
		\begin{itemize}
			\item Uses absolute differences (not squared)
			\item \textcolor{myblue}{Robust to outliers}
			\item Easy to communicate
			\item Straightforward average error magnitude
		\end{itemize}
		
		\vspace{0.1cm}
		\textbf{LIMITATIONS:}
		\begin{itemize}
			\item Can mask direction of bias (over vs under prediction)
			\item Doesn't heavily penalize occasional large misses
			\item May underestimate critical risk events from risk capital perspective
		\end{itemize}
		
		\vspace{0.1cm}
		\textbf{EXAMPLE (Salary Prediction):}
		\begin{itemize}
			\item Sum of absolute errors = 50.11
			\item N = 10
			\item MAE = 50.11 / 10 = 5.01
			\item Interpretation: "On average, predictions are off by about 5.01 thousand dollars"
		\end{itemize}
		
		\vspace{0.1cm}
		\textbf{WHEN TO USE:}
		\begin{itemize}
			\item When outliers are not critical
			\item When you want stable, interpretable measure
			\item When equal treatment of all errors is desired
		\end{itemize}
		
		\textcolor{myred}{\textbf{EXAM TRAP:}} MAE is \textcolor{myblue}{robust to outliers} - unlike MSE and RMSE!
	\end{frame}
	
	\begin{frame}{T11: MAE Managerial Question}
		\fontsize{6.5}{7.5}\selectfont
		\textbf{TOPIC: What Managerial Question Does MAE Answer?}
		
		\vspace{0.1cm}
		\textbf{MANAGERIAL QUESTION:}
		\begin{itemize}
			\item "On average, by how much are our predictions off, regardless of direction?"
		\end{itemize}
		
		\vspace{0.1cm}
		\textbf{MANAGERIAL FOCUS:}
		\begin{itemize}
			\item "On average, our predictions are off by MAE dollars"
			\item Ignores direction (over vs under prediction)
			\item Easy to communicate to non-technical audiences
		\end{itemize}
		
		\vspace{0.1cm}
		\textbf{WHEN TO USE:}
		\begin{itemize}
			\item When outliers are not critical
			\item When you want stable, interpretable measure
			\item When equal treatment of all errors is desired
		\end{itemize}
		
		\vspace{0.1cm}
		\textbf{COMPARISON WITH RMSE:}
		\begin{itemize}
			\item MAE = 5.01 vs RMSE = 6.57 in salary example
			\item RMSE > MAE because squaring penalizes larger errors more
			\item The gap between RMSE and MAE indicates presence of large errors
		\end{itemize}
		
		\vspace{0.1cm}
		\textbf{REGULATORY PERSPECTIVE:}
		\begin{itemize}
			\item ECB recommends reporting RMSE alongside MAE
			\item MAE alone might mask critical risks
			\item Both together provide complete picture
		\end{itemize}
		
		\textcolor{myred}{\textbf{EXAM TRAP:}} MAE answers \textcolor{myblue}{"by how much are we off, regardless of direction?"}
	\end{frame}
	
	\begin{frame}{T12: Mean Absolute Percentage Error (MAPE)}
		\fontsize{6.5}{7.5}\selectfont
		\textbf{TOPIC: Mean Absolute Percentage Error - Definition and Advantages}
		
		\vspace{0.1cm}
		\textbf{FORMULA:}
		\[
		MAPE = \frac{100}{N}\sum_{i=1}^N \left|\frac{y_i - \hat{y}_i}{y_i}\right|
		\]
		
		\vspace{0.1cm}
		\textbf{KEY CHARACTERISTICS:}
		\begin{itemize}
			\item Expressed as percentage
			\item Normalized - allows comparison across different scales
			\item Board-friendly due to percentage nature
			\item Consistent across portfolios
		\end{itemize}
		
		\vspace{0.1cm}
		\textbf{LIMITATIONS:}
		\begin{itemize}
			\item Explodes or becomes undefined when actual values are zero or near zero
			\item Can be misleading if few small-value items have very high percentage errors
			\item May understate errors on large absolute value items
		\end{itemize}
		
		\vspace{0.1cm}
		\textbf{EXAMPLE (Salary Prediction):}
		\begin{itemize}
			\item Sum of $|(y_i - \hat{y}_i)/y_i|$ = 1.62
			\item N = 10
			\item MAPE = 1.62 / 10 × 100 = 16.2\%
			\item Interpretation: "On average, forecast error is 16\% of actual salary"
		\end{itemize}
		
		\vspace{0.1cm}
		\textbf{MOST INTUITIVE:}
		\begin{itemize}
			\item Of all forecast accuracy measures, MAPE is most intuitive
			\item 16\% can be interpreted as average forecast error is 16\% of actual salary
		\end{itemize}
		
		\textcolor{myred}{\textbf{EXAM TRAP:}} MAPE is \textcolor{myblue}{percentage-based} - great for communication but has limitations!
	\end{frame}
	
	\begin{frame}{T13: MAPE Managerial Question}
		\fontsize{6.5}{7.5}\selectfont
		\textbf{TOPIC: What Managerial Question Does MAPE Answer?}
		
		\vspace{0.1cm}
		\textbf{MANAGERIAL QUESTION:}
		\begin{itemize}
			\item "In percentage terms, what is the typical error relative to the actual value?"
		\end{itemize}
		
		\vspace{0.1cm}
		\textbf{MANAGERIAL FOCUS:}
		\begin{itemize}
			\item "On average, our predictions are off by MAPE\%"
			\item Normalized - allows comparison across models/portfolios
			\item Perfect for audit committees and board presentations
			\item Common choice for auditors
		\end{itemize}
		
		\vspace{0.1cm}
		\textbf{WHEN TO USE:}
		\begin{itemize}
			\item Comparing models with different scales
			\item Communicating to non-technical stakeholders
			\item When percentage errors are more meaningful than absolute
		\end{itemize}
		
		\vspace{0.1cm}
		\textbf{AUDITOR PREFERENCES:}
		\begin{itemize}
			\item Auditors often gravitate to MAPE for pure percentage narrative
			\item Standardized measure for discussing model accuracy with audit committees
			\item Allows comparison across different models and portfolios
		\end{itemize}
		
		\textcolor{myred}{\textbf{EXAM TRAP:}} MAPE answers \textcolor{myblue}{"what's the error as a percentage?"}
	\end{frame}
	
	\begin{frame}{T14: MSE Calculation Example}
		\fontsize{6.5}{7.5}\selectfont
		\textbf{TOPIC: Step-by-Step MSE Calculation}
		
		\vspace{0.1cm}
		\textbf{EXAMPLE DATA (Test Sample of 10 Observations):}
		
		\begin{tabular}{|c|c|c|c|c|c|}
			\hline
			$i$ & $x_i$ & $y_i$ & $\hat{y}_i$ & $y_i - \hat{y}_i$ & $(y_i - \hat{y}_i)^2$ \\
			\hline
			1 & 4 & 20.15 & 21.12 & -0.97 & 0.94 \\
			2 & 6 & 21.48 & 23.24 & -1.76 & 3.10 \\
			3 & 9 & 31.88 & 26.42 & 5.46 & 29.81 \\
			4 & 10 & 32.88 & 27.48 & 5.40 & 29.16 \\
			5 & 21 & 49.92 & 39.14 & 10.78 & 116.21 \\
			6 & 25 & 57.46 & 43.38 & 14.08 & 198.25 \\
			7 & 3 & 21.13 & 20.06 & 1.07 & 1.14 \\
			8 & 1 & 12.18 & 17.94 & -5.76 & 33.18 \\
			9 & 2 & 19.47 & 19.00 & 0.47 & 0.22 \\
			10 & 7 & 28.66 & 24.30 & 4.36 & 19.01 \\
			\hline
			Total & & & & & 431.02 \\
			\hline
		\end{tabular}
		
		\vspace{0.1cm}
		\textbf{CALCULATION:}
		\begin{itemize}
			\item Sum of squared errors = 431.02
			\item N = 10
			\item MSE = 431.02 / 10 = 43.1
			\item RMSE = $\sqrt{43.1} = 6.57$
		\end{itemize}
		
		\textcolor{myred}{\textbf{EXAM TRAP:}} MSE = \textcolor{myblue}{sum of squared errors divided by N}!
	\end{frame}
	
	\begin{frame}{T15: RMSE Interpretation}
		\fontsize{6.5}{7.5}\selectfont
		\textbf{TOPIC: Interpreting RMSE in Context}
		
		\vspace{0.1cm}
		\textbf{EXAMPLE:}
		\begin{itemize}
			\item Salary prediction model
			\item RMSE = 6.57
			\item Target variable measured in thousands of dollars
		\end{itemize}
		
		\vspace{0.1cm}
		\textbf{INTERPRETATION:}
		\begin{itemize}
			\item "On average, predictions are off by about 6.57 thousand dollars"
			\item RMSE is in same units as target (salary in \$000s)
			\item Gives sense of typical error magnitude
		\end{itemize}
		
		\vspace{0.1cm}
		\textbf{IMPORTANT NUANCES:}
		\begin{itemize}
			\item This is an average - some errors are larger, some smaller
			\item Squaring means larger errors contribute more to the average
			\item RMSE = 6.57 means typical error is around 6.57 units
			\item Not every prediction is off by exactly 6.57
		\end{itemize}
		
		\vspace{0.1cm}
		\textbf{COMPARISON WITH MAE:}
		\begin{itemize}
			\item MAE = 5.01
			\item RMSE = 6.57
			\item Gap between RMSE and MAE indicates presence of larger errors
			\item If RMSE $\approx$ MAE, errors are relatively uniform
		\end{itemize}
		
		\textcolor{myred}{\textbf{EXAM TRAP:}} RMSE is in \textcolor{myblue}{same units as target} - easier to understand!
	\end{frame}
	
	% ============================================================
	% SECTION 3: METRIC COMPARISON AND SELECTION
	% ============================================================
	
	\begin{frame}{T16: MAE Calculation Example}
		\fontsize{6.5}{7.5}\selectfont
		\textbf{TOPIC: Step-by-Step MAE Calculation}
		
		\vspace{0.1cm}
		\textbf{EXAMPLE DATA (Test Sample of 10 Observations):}
		
		\begin{tabular}{|c|c|c|c|c|c|}
			\hline
			$i$ & $x_i$ & $y_i$ & $\hat{y}_i$ & $y_i - \hat{y}_i$ & $|y_i - \hat{y}_i|$ \\
			\hline
			1 & 4 & 20.15 & 21.12 & -0.97 & 0.97 \\
			2 & 6 & 21.48 & 23.24 & -1.76 & 1.76 \\
			3 & 9 & 31.88 & 26.42 & 5.46 & 5.46 \\
			4 & 10 & 32.88 & 27.48 & 5.40 & 5.40 \\
			5 & 21 & 49.92 & 39.14 & 10.78 & 10.78 \\
			6 & 25 & 57.46 & 43.38 & 14.08 & 14.08 \\
			7 & 3 & 21.13 & 20.06 & 1.07 & 1.07 \\
			8 & 1 & 12.18 & 17.94 & -5.76 & 5.76 \\
			9 & 2 & 19.47 & 19.00 & 0.47 & 0.47 \\
			10 & 7 & 28.66 & 24.30 & 4.36 & 4.36 \\
			\hline
			Total & & & & & 50.11 \\
			\hline
		\end{tabular}
		
		\vspace{0.1cm}
		\textbf{CALCULATION:}
		\begin{itemize}
			\item Sum of absolute errors = 50.11
			\item N = 10
			\item MAE = 50.11 / 10 = 5.01
		\end{itemize}
		
		\vspace{0.1cm}
		\textbf{INTERPRETATION:}
		\begin{itemize}
			\item "On average, predictions are off by about 5.01 thousand dollars"
			\item Similar to RMSE but less influenced by large errors
			\item More stable measure of typical error
		\end{itemize}
		
		\textcolor{myred}{\textbf{EXAM TRAP:}} MAE = \textcolor{myblue}{average of absolute errors} - sum divided by N!
	\end{frame}
	
	\begin{frame}{T17: MAPE Calculation Example}
		\fontsize{6.5}{7.5}\selectfont
		\textbf{TOPIC: Step-by-Step MAPE Calculation}
		
		\vspace{0.1cm}
		\textbf{EXAMPLE DATA (Test Sample of 10 Observations):}
		
		\begin{tabular}{|c|c|c|c|c|c|}
			\hline
			$i$ & $y_i$ & $\hat{y}_i$ & $y_i - \hat{y}_i$ & $\frac{y_i - \hat{y}_i}{y_i}$ & $\left|\frac{y_i - \hat{y}_i}{y_i}\right|$ \\
			\hline
			1 & 20.15 & 21.12 & -0.97 & -0.05 & 0.05 \\
			2 & 21.48 & 23.24 & -1.76 & -0.08 & 0.08 \\
			3 & 31.88 & 26.42 & 5.46 & 0.17 & 0.17 \\
			4 & 32.88 & 27.48 & 5.40 & 0.16 & 0.16 \\
			5 & 49.92 & 39.14 & 10.78 & 0.22 & 0.22 \\
			6 & 57.46 & 43.38 & 14.08 & 0.25 & 0.25 \\
			7 & 21.13 & 20.06 & 1.07 & 0.05 & 0.05 \\
			8 & 12.18 & 17.94 & -5.76 & -0.47 & 0.47 \\
			9 & 19.47 & 19.00 & 0.47 & 0.02 & 0.02 \\
			10 & 28.66 & 24.30 & 4.36 & 0.15 & 0.15 \\
			\hline
			Total & & & & & 1.62 \\
			\hline
		\end{tabular}
		
		\vspace{0.1cm}
		\textbf{CALCULATION:}
		\begin{itemize}
			\item Sum of $|(y_i - \hat{y}_i)/y_i|$ = 1.62
			\item N = 10
			\item MAPE = 1.62 / 10 × 100 = 16.2\%
		\end{itemize}
		
		\vspace{0.1cm}
		\textbf{INTERPRETATION:}
		\begin{itemize}
			\item "On average, forecast error is 16\% of actual salary"
			\item Most intuitive of all measures
			\item Easy to communicate to non-technical stakeholders
		\end{itemize}
		
		\textcolor{myred}{\textbf{EXAM TRAP:}} MAPE = \textcolor{myblue}{percentage of actual values} - most intuitive!
	\end{frame}
	
	\begin{frame}{T18: Metric Selection for Tail Risk}
		\fontsize{6.5}{7.5}\selectfont
		\textbf{TOPIC: Choosing Metrics to Capture Tail Risk}
		
		\vspace{0.1cm}
		\textbf{WHICH METRIC CAPTURES TAIL RISK?}
		\begin{itemize}
			\item \textcolor{myblue}{MSE or RMSE} - because they square errors, heavily penalizing large deviations
		\end{itemize}
		
		\vspace{0.1cm}
		\textbf{WHY MSE/RMSE CAPTURE TAIL RISK:}
		\begin{itemize}
			\item Squaring errors heavily penalizes large deviations
			\item Large errors contribute disproportionately to the metric
			\item Sensitive to outliers (tail events)
			\item Regulatory preference (ECB recommends)
		\end{itemize}
		
		\vspace{0.1cm}
		\textbf{COMPARISON:}
		\begin{itemize}
			\item \textcolor{myblue}{MAE:} Robust, but ignores tail risk
			\item \textcolor{myblue}{MAPE:} Good for communication, but has limitations
			\item \textcolor{myblue}{MSE/RMSE:} Best for capturing extreme errors
		\end{itemize}
		
		\vspace{0.1cm}
		\textbf{REGULATORY PERSPECTIVE:}
		\begin{itemize}
			\item ECB's Guide to Internal Models recommends metrics capturing tail risk
			\item Nudges firms to report RMSE alongside MAE
			\item Ensures large, infrequent losses are not overlooked
		\end{itemize}
		
		\textcolor{myred}{\textbf{EXAM TRAP:}} ECB recommends metrics that \textcolor{myblue}{capture tail risk} - MSE and RMSE!
	\end{frame}
	
	\begin{frame}{T19: When to Use MAPE}
		\fontsize{6.5}{7.5}\selectfont
		\textbf{TOPIC: Appropriate Use Cases for MAPE}
		
		\vspace{0.1cm}
		\textbf{WHEN MAPE IS APPROPRIATE:}
		\begin{itemize}
			\item Comparing models across different portfolios
			\item Comparing error on €100k loan vs €10m loan
			\item Communicating to audit committees
			\item When percentage errors are more meaningful than absolute
		\end{itemize}
		
		\vspace{0.1cm}
		\textbf{WHEN MAPE IS NOT APPROPRIATE:}
		\begin{itemize}
			\item When actual values are near zero (MAPE explodes)
			\item When tail risk is the primary concern
			\item When absolute errors matter more than percentage
		\end{itemize}
		
		\vspace{0.1cm}
		\textbf{KEY ADVANTAGE:}
		\begin{itemize}
			\item Normalized measure allows comparison across scales
			\item A €2,000 error on €10,000 loan (20\%) vs €2,000 error on €1,000,000 loan (0.2\%)
			\item MAPE treats these appropriately based on relative size
		\end{itemize}
		
		\textcolor{myred}{\textbf{EXAM TRAP:}} MAPE is best for \textcolor{myblue}{comparing across different scales}!
	\end{frame}
	
	\begin{frame}{T20: Error Distribution Visualization}
		\fontsize{6.5}{7.5}\selectfont
		\textbf{TOPIC: Beyond Averages - The Story in Error Distributions}
		
		\vspace{0.1cm}
		\textbf{WHY LOOK BEYOND AVERAGES?}
		\begin{itemize}
			\item Distribution of errors reveals patterns averages hide
			\item Senior leaders should encourage visualizations (histograms, box plots)
			\item Provides richer context for interpreting average performance metrics
		\end{itemize}
		
		\vspace{0.1cm}
		\textbf{KEY QUESTIONS TO ASK:}
		\begin{enumerate}
			\item "Are our model errors generally small with a few large ones, or consistently off by moderate amount?"
			\item "Are errors randomly distributed or is model systematically over-predicting (more negative errors) or under-predicting (more positive errors)?"
			\item "Are there particular segments or conditions where errors are consistently larger or more biased?"
		\end{enumerate}
		
		\vspace{0.1cm}
		\textbf{WHY THIS MATTERS:}
		\begin{itemize}
			\item \textcolor{myblue}{Consistent bias} > random errors (more problematic even if average small)
			\item Segment analysis reveals model weaknesses
			\item Richer context for interpreting average metrics
			\item Diagnoses underlying model problems
		\end{itemize}
		
		\textcolor{myred}{\textbf{EXAM TRAP:}} Error distributions reveal \textcolor{myblue}{patterns and biases} that averages hide!
	\end{frame}
	
	\begin{frame}{T21: Model Comparison Example}
		\fontsize{6.5}{7.5}\selectfont
		\textbf{TOPIC: Comparing Alternative Models}
		
		\vspace{0.1cm}
		\textbf{EXAMPLE: House Price Prediction}
		
		\begin{tabular}{|p{3.5cm}|p{2.5cm}|p{2.5cm}|p{2.5cm}|}
			\hline
			\textbf{Model} & \textbf{MSE} & \textbf{MAE} & \textbf{MAPE} \\
			\hline
			Model A (Linear Regression) & 94.85 & 6.56 & 18.25 \\
			\hline
			Model B (Decision Tree) & 77.35* & 5.41* & 14.45* \\
			\hline
		\end{tabular}
		
		\vspace{0.1cm}
		\textbf{CONCLUSION:}
		\begin{itemize}
			\item Model B has all lower values
			\item Decision tree is more accurate according to all metrics
			\item Asterisk denotes best performer
		\end{itemize}
		
		\vspace{0.1cm}
		\textbf{KEY INSIGHT:}
		\begin{itemize}
			\item Different metrics can yield different model rankings
			\item Choice depends on problem, regulatory situation, and historical precedent
			\item Lower values = better performance for all these metrics
		\end{itemize}
		
		\textcolor{myred}{\textbf{EXAM TRAP:}} Lower values = \textcolor{myblue}{better performance} for all these metrics!
	\end{frame}
	
	\begin{frame}{T22: Different Metrics Different Rankings}
		\fontsize{6.5}{7.5}\selectfont
		\textbf{TOPIC: Why Different Metrics Can Yield Different Rankings}
		
		\vspace{0.1cm}
		\textbf{WHY RANKINGS CAN DIFFER:}
		\begin{itemize}
			\item \textcolor{myblue}{MSE/RMSE:} Penalize large errors heavily
			\item \textcolor{myblue}{MAE:} Treats all errors equally
			\item \textcolor{myblue}{MAPE:} Emphasizes relative errors
			\item Different metrics = different priorities
		\end{itemize}
		
		\vspace{0.1cm}
		\textbf{IMPLICATION:}
		\begin{itemize}
			\item Choose metric based on business needs
			\item Regulatory requirements may dictate choice
			\item Historical organizational choices matter
			\item No single "best" metric for all situations
		\end{itemize}
		
		\vspace{0.1cm}
		\textbf{EXAMPLE SCENARIO:}
		\begin{itemize}
			\item Model A: Lower MAE but higher RMSE
			\item Model B: Lower RMSE but higher MAE
			\item Choice depends on whether tail risk or average error matters more
		\end{itemize}
		
		\textcolor{myred}{\textbf{EXAM TRAP:}} Different metrics = \textcolor{myblue}{different rankings} - choose based on problem!
	\end{frame}
	
	\begin{frame}{T23: Unnormalized Metrics Limitation}
		\fontsize{6.5}{7.5}\selectfont
		\textbf{TOPIC: Limitations of Unnormalized Metrics}
		
		\vspace{0.1cm}
		\textbf{THE PROBLEM:}
		\begin{itemize}
			\item RMSE and MAE are both unnormalized measures
			\item They treat errors symmetrically regardless of scale
			\item Example: \$2,000 error on \$150,000 salary (1.3\%) vs \$2,000 error on \$12,000 salary (16.7\%)
			\item Both contribute same to RMSE/MAE
			\item But the second error is much more significant proportionally
		\end{itemize}
		
		\vspace{0.1cm}
		\textbf{WHY THIS MATTERS:}
		\begin{itemize}
			\item In some situations, a \$2,000 error on \$150,000 might be considered much less important
			\item Represents about 1\% of total salary
			\item The other situation represents more than 10\% of salary
			\item Unnormalized metrics don't capture this distinction
		\end{itemize}
		
		\vspace{0.1cm}
		\textbf{SOLUTION:}
		\begin{itemize}
			\item Use MAPE for relative comparison
			\item Understand the context of errors
			\item Consider business impact, not just statistical measures
		\end{itemize}
		
		\textcolor{myred}{\textbf{EXAM TRAP:}} Unnormalized metrics treat \textcolor{myblue}{all errors equally} - MAPE handles this issue!
	\end{frame}
	
	\begin{frame}{T24: Regulatory Metric Preferences}
		\fontsize{6.5}{7.5}\selectfont
		\textbf{TOPIC: ECB Recommendations on Performance Metrics}
		
		\vspace{0.1cm}
		\textbf{ECB GUIDE TO INTERNAL MODELS:}
		\begin{itemize}
			\item Recommends use of metrics that capture tail risk
			\item Nudges firms to report RMSE alongside MAE
			\item Ensures that impact of large, infrequent losses is not overlooked
		\end{itemize}
		
		\vspace{0.1cm}
		\textbf{WHY THIS MATTERS:}
		\begin{itemize}
			\item MAE alone might mask critical risks
			\item RMSE captures the impact of extreme events
			\item Both together provide complete picture
			\item Aligns with regulatory capital frameworks
		\end{itemize}
		
		\vspace{0.1cm}
		\textbf{REGULATORY CONTEXT:}
		\begin{itemize}
			\item ECB's Guide to Internal Models
			\item Focus on tail risk and extreme events
			\item Capital adequacy requirements
			\item Systemic risk concerns
		\end{itemize}
		
		\textcolor{myred}{\textbf{EXAM TRAP:}} ECB recommends \textcolor{myblue}{RMSE alongside MAE} for complete risk picture!
	\end{frame}
	
	\begin{frame}{T25: Auditor Preferences}
		\fontsize{6.5}{7.5}\selectfont
		\textbf{TOPIC: Why Auditors Prefer MAPE}
		
		\vspace{0.1cm}
		\textbf{AUDITOR PERSPECTIVE:}
		\begin{itemize}
			\item Auditors often gravitate to MAPE
			\item Pure percentage narrative when discussing model accuracy with audit committees
			\item Provides standardized measure
			\item Allows comparison across different models/portfolios
		\end{itemize}
		
		\vspace{0.1cm}
		\textbf{WHY PERCENTAGE NARRATIVE MATTERS:}
		\begin{itemize}
			\item Easy to understand for non-technical stakeholders
			\item Standardized across different scales
			\item Facilitates comparison
			\item Communicates effectively with audit committees
		\end{itemize}
		
		\vspace{0.1cm}
		\textbf{DIFFERENT STAKEHOLDER PREFERENCES:}
		\begin{itemize}
			\item \textcolor{myblue}{Regulators:} RMSE alongside MAE (tail risk)
			\item \textcolor{myblue}{Auditors:} MAPE (percentage narrative)
			\item \textcolor{myblue}{Risk Managers:} All metrics for complete picture
			\item \textcolor{myblue}{Board:} MAPE and accuracy (intuitive)
		\end{itemize}
		
		\textcolor{myred}{\textbf{EXAM TRAP:}} Auditors prefer MAPE for its \textcolor{myblue}{percentage narrative}!
	\end{frame}
	
	% ============================================================
	% SECTION 4: CONFUSION MATRIX AND CLASSIFICATION METRICS
	% ============================================================
	
	\begin{frame}{T26: Confusion Matrix Definition}
		\fontsize{6.5}{7.5}\selectfont
		\textbf{TOPIC: The Confusion Matrix - Foundation of Classification Metrics}
		
		\vspace{0.1cm}
		\textbf{DEFINITION:}
		\begin{itemize}
			\item Cross-tabulation of predicted classifications against actual outcomes
			\item Simple yet powerful tool providing clear breakdown of model performance
			\item Cornerstone for evaluating classification models
		\end{itemize}
		
		\vspace{0.1cm}
		\textbf{STRUCTURE:}
		
		\begin{tabular}{|c|c|c|}
			\hline
			& \textbf{Predicted Negative} & \textbf{Predicted Positive} \\
			\hline
			\textbf{Actual Negative} & True Negative (TN) & False Positive (FP) \\
			\hline
			\textbf{Actual Positive} & False Negative (FN) & True Positive (TP) \\
			\hline
		\end{tabular}
		
		\vspace{0.1cm}
		\textbf{DEFINITIONS FROM RISK MANAGER'S PERSPECTIVE:}
		\begin{itemize}
			\item \textcolor{myblue}{True Negative (TN):} Model correctly identified negative case
			\item \textcolor{myblue}{False Positive (FP):} Model incorrectly classified negative as positive (Type I Error)
			\item \textcolor{myblue}{False Negative (FN):} Model incorrectly classified positive as negative (Type II Error)
			\item \textcolor{myblue}{True Positive (TP):} Model correctly identified positive case
		\end{itemize}
		
		\textcolor{myred}{\textbf{EXAM TRAP:}} Confusion matrix = \textcolor{myblue}{predicted vs actual} cross-tabulation!
	\end{frame}
	
	\begin{frame}{T27: True Positive (TP)}
		\fontsize{6.5}{7.5}\selectfont
		\textbf{TOPIC: True Positive - Correctly Identified Positive Case}
		
		\vspace{0.1cm}
		\textbf{DEFINITION:}
		\begin{itemize}
			\item The model correctly identified a positive case
			\item Desired outcome - what we want the model to do
		\end{itemize}
		
		\vspace{0.1cm}
		\textbf{EXAMPLES:}
		\begin{itemize}
			\item Correctly flagged a fraudulent transaction
			\item Correctly predicted that a loan would default
			\item Correctly identified a dividend-paying firm
			\item Correctly detected a critical compliance breach
		\end{itemize}
		
		\vspace{0.1cm}
		\textbf{WHY IT MATTERS:}
		\begin{itemize}
			\item More TP = better model performance
			\item True positives represent value creation
			\item Catching fraud, defaults, or breaches avoids losses
		\end{itemize}
		
		\vspace{0.1cm}
		\textbf{EXAMPLE (Dividend Prediction):}
		\begin{itemize}
			\item TP = 432
			\item 432 firms correctly predicted to pay dividends
			\item These are the successful predictions
		\end{itemize}
		
		\textcolor{myred}{\textbf{EXAM TRAP:}} True Positive = \textcolor{myblue}{correctly identified positive}!
	\end{frame}
	
	\begin{frame}{T28: True Negative (TN)}
		\fontsize{6.5}{7.5}\selectfont
		\textbf{TOPIC: True Negative - Correctly Identified Negative Case}
		
		\vspace{0.1cm}
		\textbf{DEFINITION:}
		\begin{itemize}
			\item The model correctly identified a negative case
			\item Desired outcome - what we want the model to do
		\end{itemize}
		
		\vspace{0.1cm}
		\textbf{EXAMPLES:}
		\begin{itemize}
			\item Correctly identified a legitimate transaction
			\item Correctly predicted that a loan would not default
			\item Correctly identified a non-dividend paying firm
			\item Correctly cleared a compliant process
		\end{itemize}
		
		\vspace{0.1cm}
		\textbf{WHY IT MATTERS:}
		\begin{itemize}
			\item More TN = better model performance
			\item True negatives represent efficiency
			\item Avoiding unnecessary interventions or investigations
		\end{itemize}
		
		\vspace{0.1cm}
		\textbf{EXAMPLE (Dividend Prediction):}
		\begin{itemize}
			\item TN = 279
			\item 279 firms correctly predicted to not pay dividends
			\item These are the successful negative predictions
		\end{itemize}
		
		\textcolor{myred}{\textbf{EXAM TRAP:}} True Negative = \textcolor{myblue}{correctly identified negative}!
	\end{frame}
	
	\begin{frame}{T29: False Positive (FP)}
		\fontsize{6.5}{7.5}\selectfont
		\textbf{TOPIC: False Positive - Type I Error}
		
		\vspace{0.1cm}
		\textbf{DEFINITION:}
		\begin{itemize}
			\item The model incorrectly classified a negative case as positive
			\item Also known as Type I Error
			\item "False alarm"
		\end{itemize}
		
		\vspace{0.1cm}
		\textbf{EXAMPLES:}
		\begin{itemize}
			\item Flagged a legitimate transaction as fraudulent
			\item Predicted that a good loan would default
			\item Incorrectly identified a firm as paying dividends
			\item Wrongly flagged a compliant process as breach
		\end{itemize}
		
		\vspace{0.1cm}
		\textbf{COST PERSPECTIVE:}
		\begin{itemize}
			\item Operational cost to investigate false alert
			\item Customer friction from declined transactions
			\item Lost revenue from declining good customers
			\item Resource waste from chasing red herrings
		\end{itemize}
		
		\vspace{0.1cm}
		\textbf{EXAMPLE (Dividend Prediction):}
		\begin{itemize}
			\item FP = 121
			\item 121 firms incorrectly predicted to pay dividends
			\item False alarms in prediction
		\end{itemize}
		
		\textcolor{myred}{\textbf{EXAM TRAP:}} False Positive = \textcolor{myblue}{Type I Error} - incorrectly saying "yes"!
	\end{frame}
	
	\begin{frame}{T30: False Negative (FN)}
		\fontsize{6.5}{7.5}\selectfont
		\textbf{TOPIC: False Negative - Type II Error}
		
		\vspace{0.1cm}
		\textbf{DEFINITION:}
		\begin{itemize}
			\item The model incorrectly classified a positive case as negative
			\item Also known as Type II Error
			\item "Missed detection"
		\end{itemize}
		
		\vspace{0.1cm}
		\textbf{EXAMPLES:}
		\begin{itemize}
			\item Missed a fraudulent transaction (labeled it legitimate)
			\item Failed to predict that a loan would default
			\item Incorrectly identified a firm as not paying dividends
			\item Missed a critical compliance breach
		\end{itemize}
		
		\vspace{0.1cm}
		\textbf{COST PERSPECTIVE:}
		\begin{itemize}
			\item Often the most costly type of error
			\item Financial losses from undetected fraud
			\item Capital impact of missed defaults
			\item Regulatory sanctions for missed breaches
		\end{itemize}
		
		\vspace{0.1cm}
		\textbf{EXAMPLE (Dividend Prediction):}
		\begin{itemize}
			\item FN = 168
			\item 168 firms incorrectly predicted to not pay dividends
			\item Missed dividend payers
		\end{itemize}
		
		\textcolor{myred}{\textbf{EXAM TRAP:}} False Negative = \textcolor{myblue}{Type II Error} - incorrectly saying "no"!
	\end{frame}
	
	\begin{frame}{T31: Type I vs Type II Errors}
		\fontsize{6.5}{7.5}\selectfont
		\textbf{TOPIC: Comparing Type I and Type II Errors in Credit Risk}
		
		\vspace{0.1cm}
		\textbf{WHY TYPE II ERRORS ARE MORE COSTLY IN CREDIT RISK:}
		\begin{itemize}
			\item Misclassifying a borrower as solvent who then defaults
			\item Results in actual loan losses for the bank
			\item Regulatory capital implications
			\item Direct financial impact
		\end{itemize}
		
		\vspace{0.1cm}
		\textbf{TYPE I ERROR COST:}
		\begin{itemize}
			\item Lost profit from not making a loan
			\item Operational cost of investigation
			\item Customer friction
			\item Usually less severe than Type II
		\end{itemize}
		
		\vspace{0.1cm}
		\textbf{GENERAL PRINCIPLE:}
		\begin{itemize}
			\item The cost of errors is typically not symmetric
			\item Choose threshold based on relative costs
			\item In credit risk, Type II errors usually more costly
			\item In AML, both types have significant costs
		\end{itemize}
		
		\vspace{0.1cm}
		\textbf{EXAMPLES:}
		\begin{itemize}
			\item \textcolor{myblue}{Credit Risk:} Type II more costly (missed default)
			\item \textcolor{myblue}{Fraud Detection:} Type II more costly (missed fraud)
			\item \textcolor{myblue}{AML:} Both costly (missed suspicious activity vs customer friction)
		\end{itemize}
		
		\textcolor{myred}{\textbf{EXAM TRAP:}} Type II Errors (False Negatives) are \textcolor{myblue}{usually more costly} in credit risk!
	\end{frame}
	
	\begin{frame}{T32: Accuracy Definition}
		\fontsize{6.5}{7.5}\selectfont
		\textbf{TOPIC: Accuracy - Basic Classification Performance Measure}
		
		\vspace{0.1cm}
		\textbf{FORMULA:}
		\[
		Accuracy = \frac{TP + TN}{TP + TN + FP + FN}
		\]
		
		\vspace{0.1cm}
		\textbf{INTERPRETATION:}
		\begin{itemize}
			\item "Percentage of all predictions that were correct"
			\item Sum of diagonal elements of confusion matrix divided by total
			\item Straightforward to interpret - simple agreement between predicted and observed
		\end{itemize}
		
		\vspace{0.1cm}
		\textbf{EXAMPLE (Dividend Prediction):}
		\begin{itemize}
			\item Accuracy = (432 + 279) / 1000 = 71.1\%
			\item "71.1\% of dividend predictions were correct"
			\item Error rate = 100\% - 71.1\% = 28.9\%
		\end{itemize}
		
		\vspace{0.1cm}
		\textbf{LIMITATIONS:}
		\begin{itemize}
			\item Doesn't distinguish between error types (FP and FN treated equally)
			\item Costs of errors are typically not symmetric
			\item Misleading for imbalanced datasets
			\item 98\% positive, classify all as positive → 98\% accuracy but useless model
		\end{itemize}
		
		\textcolor{myred}{\textbf{EXAM TRAP:}} Accuracy = \textcolor{myblue}{percentage of all correct predictions}!
	\end{frame}
	
	\begin{frame}{T33: Accuracy Example}
		\fontsize{6.5}{7.5}\selectfont
		\textbf{TOPIC: Calculating Accuracy from Confusion Matrix}
		
		\vspace{0.1cm}
		\textbf{EXAMPLE DATA:}
		\begin{itemize}
			\item True Positive (TP) = 432
			\item True Negative (TN) = 279
			\item False Positive (FP) = 121
			\item False Negative (FN) = 168
		\end{itemize}
		
		\vspace{0.1cm}
		\textbf{CALCULATION:}
		\begin{itemize}
			\item Accuracy = (TP + TN) / (TP + TN + FP + FN)
			\item Accuracy = (432 + 279) / (432 + 279 + 121 + 168)
			\item Accuracy = 711 / 1000
			\item Accuracy = 71.1\%
		\end{itemize}
		
		\vspace{0.1cm}
		\textbf{INTERPRETATION:}
		\begin{itemize}
			\item "71.1\% of all predictions were correct"
			\item Error rate = 28.9\% (percentage of incorrect predictions)
			\item 711 out of 1000 predictions were correct
		\end{itemize}
		
		\vspace{0.1cm}
		\textbf{CONTEXT:}
		\begin{itemize}
			\item 600 firms actually paid dividends
			\item 400 firms actually did not
			\item Model correctly classified 71.1\% of all firms
		\end{itemize}
		
		\textcolor{myred}{\textbf{EXAM TRAP:}} Accuracy = \textcolor{myblue}{(TP + TN) / Total}!
	\end{frame}
	
	\begin{frame}{T34: Error Rate Definition}
		\fontsize{6.5}{7.5}\selectfont
		\textbf{TOPIC: Error Rate and Its Relationship to Accuracy}
		
		\vspace{0.1cm}
		\textbf{FORMULA:}
		\[
		Error Rate = 1 - Accuracy = \frac{FP + FN}{TP + TN + FP + FN}
		\]
		
		\vspace{0.1cm}
		\textbf{INTERPRETATION:}
		\begin{itemize}
			\item "Percentage of predictions that were incorrect"
			\item Complementary to accuracy
			\item Simple and intuitive
		\end{itemize}
		
		\vspace{0.1cm}
		\textbf{EXAMPLE:}
		\begin{itemize}
			\item Accuracy = 71.1\%
			\item Error Rate = 100\% - 71.1\% = 28.9\%
			\item "28.9\% of predictions were wrong"
		\end{itemize}
		
		\vspace{0.1cm}
		\textbf{RELATIONSHIP:}
		\begin{itemize}
			\item Error Rate = 1 - Accuracy
			\item They are complementary measures
			\item Both ignore the type of error made
			\item Both problematic for imbalanced datasets
		\end{itemize}
		
		\textcolor{myred}{\textbf{EXAM TRAP:}} Error rate = \textcolor{myblue}{1 - Accuracy}!
	\end{frame}
	
	\begin{frame}{T35: Accuracy Limitations}
		\fontsize{6.5}{7.5}\selectfont
		\textbf{TOPIC: Major Limitations of Accuracy}
		
		\vspace{0.1cm}
		\textbf{LIMITATION 1: DOESN'T DISTINGUISH ERROR TYPES}
		\begin{itemize}
			\item FP and FN treated equally
			\item But costs are usually different
			\item Type II errors often more costly in credit risk
		\end{itemize}
		
		\vspace{0.1cm}
		\textbf{LIMITATION 2: MISLEADING FOR IMBALANCED DATA}
		\begin{itemize}
			\item 98\% positive, classify all as positive → 98\% accuracy
			\item But model is practically useless
			\item Example: Fraud detection - 99\% legitimate
			\item Model predicting "legitimate" always has 99\% accuracy
			\item But catches zero fraud (Recall = 0)
		\end{itemize}
		
		\vspace{0.1cm}
		\textbf{LIMITATION 3: IGNORES COSTS}
		\begin{itemize}
			\item Doesn't account for different costs of FP vs FN
			\item In credit risk, FN (missed default) much more costly
		\end{itemize}
		
		\vspace{0.1cm}
		\textbf{WHAT TO USE INSTEAD:}
		\begin{itemize}
			\item Precision, Recall, F1 Score
			\item AUC (Area Under ROC Curve)
			\item Cost-sensitive metrics
		\end{itemize}
		
		\textcolor{myred}{\textbf{EXAM TRAP:}} Accuracy is \textcolor{myblue}{misleading for imbalanced datasets}!
	\end{frame}
	
	\begin{frame}{T36: Precision Definition}
		\fontsize{6.5}{7.5}\selectfont
		\textbf{TOPIC: Precision - Of Predicted Positives, How Many Were Correct?}
		
		\vspace{0.1cm}
		\textbf{FORMULA:}
		\[
		Precision = \frac{TP}{TP + FP}
		\]
		
		\vspace{0.1cm}
		\textbf{MANAGERIAL QUESTION:}
		\begin{itemize}
			\item "Of all the alerts our model generated, how many were actually real issues?"
			\item "If we act on this prediction, how often will we be right?"
		\end{itemize}
		
		\vspace{0.1cm}
		\textbf{INTERPRETATION:}
		\begin{itemize}
			\item Estimate of probability that model is right when labelling an outcome as true
			\item Of all instances classified as positive, what fraction were actually positive?
		\end{itemize}
		
		\vspace{0.1cm}
		\textbf{WHEN PRECISION MATTERS:}
		\begin{itemize}
			\item High cost of investigation (AML alerts)
			\item Customer friction concerns
			\item Operational resource constraints
			\item When false alarms are expensive
		\end{itemize}
		
		\vspace{0.1cm}
		\textbf{EXAMPLE (Dividend Prediction):}
		\begin{itemize}
			\item Precision = 432 / (432 + 121) = 432/553 = 78.1\%
			\item "When model predicts dividend, it's right 78.1\% of the time"
		\end{itemize}
		
		\textcolor{myred}{\textbf{EXAM TRAP:}} Precision = \textcolor{myblue}{TP / (TP + FP)} - "of predicted positives, how many were right?"
	\end{frame}
	
	\begin{frame}{T37: Precision Example}
		\fontsize{6.5}{7.5}\selectfont
		\textbf{TOPIC: Calculating Precision from Confusion Matrix}
		
		\vspace{0.1cm}
		\textbf{EXAMPLE DATA:}
		\begin{itemize}
			\item True Positive (TP) = 432
			\item False Positive (FP) = 121
		\end{itemize}
		
		\vspace{0.1cm}
		\textbf{CALCULATION:}
		\begin{itemize}
			\item Precision = TP / (TP + FP)
			\item Precision = 432 / (432 + 121)
			\item Precision = 432 / 553
			\item Precision = 78.1\%
		\end{itemize}
		
		\vspace{0.1cm}
		\textbf{INTERPRETATION:}
		\begin{itemize}
			\item "When the model predicts a firm will pay a dividend, it's right 78.1\% of the time"
			\item "Of all dividend predictions, 78.1\% were correct"
		\end{itemize}
		
		\vspace{0.1cm}
		\textbf{CONTEXT:}
		\begin{itemize}
			\item 432 firms correctly predicted to pay dividends
			\item 121 firms incorrectly predicted to pay dividends
			\item Total predicted positive = 553
			\item Precision = 78.1\%
		\end{itemize}
		
		\textcolor{myred}{\textbf{EXAM TRAP:}} Precision = \textcolor{myblue}{TP / (TP + FP)}!
	\end{frame}
	
	\begin{frame}{T38: Recall Definition}
		\fontsize{6.5}{7.5}\selectfont
		\textbf{TOPIC: Recall (Sensitivity) - Of Actual Positives, How Many Did We Catch?}
		
		\vspace{0.1cm}
		\textbf{FORMULA:}
		\[
		Recall = \frac{TP}{TP + FN}
		\]
		
		\vspace{0.1cm}
		\textbf{ALSO KNOWN AS:} Sensitivity, True Positive Rate
		
		\vspace{0.1cm}
		\textbf{MANAGERIAL QUESTION:}
		\begin{itemize}
			\item "Of all real fraud events that occurred, how many did we catch?"
			\item "What percentage of actual positives did we identify?"
		\end{itemize}
		
		\vspace{0.1cm}
		\textbf{INTERPRETATION:}
		\begin{itemize}
			\item Proportion of correctly classified positives among all positive instances
			\item True positive rate
		\end{itemize}
		
		\vspace{0.1cm}
		\textbf{WHEN RECALL MATTERS:}
		\begin{itemize}
			\item Missing a positive has severe consequences
			\item Fraud detection, default prediction
			\item Compliance monitoring
		\end{itemize}
		
		\vspace{0.1cm}
		\textbf{SPECIFICITY (TRUE NEGATIVE RATE):}
		\[
		Specificity = \frac{TN}{TN + FP}
		\]
		
		\vspace{0.1cm}
		\textbf{EXAMPLE (Dividend Prediction):}
		\begin{itemize}
			\item Recall = 432 / (432 + 168) = 432/600 = 72.0\%
			\item "Of all firms that paid dividends, model caught 72\%"
		\end{itemize}
		
		\textcolor{myred}{\textbf{EXAM TRAP:}} Recall = \textcolor{myblue}{TP / (TP + FN)} - "of actual positives, how many did we catch?"
	\end{frame}
	
	\begin{frame}{T39: Recall Example}
		\fontsize{6.5}{7.5}\selectfont
		\textbf{TOPIC: Calculating Recall from Confusion Matrix}
		
		\vspace{0.1cm}
		\textbf{EXAMPLE DATA:}
		\begin{itemize}
			\item True Positive (TP) = 432
			\item False Negative (FN) = 168
		\end{itemize}
		
		\vspace{0.1cm}
		\textbf{CALCULATION:}
		\begin{itemize}
			\item Recall = TP / (TP + FN)
			\item Recall = 432 / (432 + 168)
			\item Recall = 432 / 600
			\item Recall = 72.0\%
		\end{itemize}
		
		\vspace{0.1cm}
		\textbf{INTERPRETATION:}
		\begin{itemize}
			\item "Of all firms that actually paid dividends, the model caught 72\%"
			\item "The model missed 28\% of dividend payers"
		\end{itemize}
		
		\vspace{0.1cm}
		\textbf{CONTEXT:}
		\begin{itemize}
			\item 600 firms actually paid dividends
			\item 432 were correctly predicted
			\item 168 were missed (False Negatives)
			\item Recall = 72.0\%
		\end{itemize}
		
		\textcolor{myred}{\textbf{EXAM TRAP:}} Recall = \textcolor{myblue}{TP / (TP + FN)}!
	\end{frame}
	
	\begin{frame}{T40: Precision vs Recall Tradeoff}
		\fontsize{6.5}{7.5}\selectfont
		\textbf{TOPIC: The Tradeoff Between Precision and Recall}
		
		\vspace{0.1cm}
		\textbf{THE TRADEOFF:}
		\begin{itemize}
			\item \textcolor{myblue}{High threshold:} Higher precision, lower recall
			\item \textcolor{myblue}{Low threshold:} Lower precision, higher recall
			\item Cannot maximize both simultaneously
		\end{itemize}
		
		\vspace{0.1cm}
		\textbf{EXAMPLE:}
		\begin{itemize}
			\item Threshold = 0.9: Only very confident predictions
			\item Precision: High (few false positives)
			\item Recall: Low (miss many actual positives)
		\end{itemize}
		
		\vspace{0.1cm}
		\textbf{BUSINESS CONTEXT DETERMINES BALANCE:}
		\begin{itemize}
			\item \textcolor{myblue}{Fraud detection:} Prioritize recall (catch all fraud)
			\item \textcolor{myblue}{AML monitoring:} High recall critical
			\item \textcolor{myblue}{Credit approval:} Balance needed
			\item \textcolor{myblue}{Marketing:} Precision more important
		\end{itemize}
		
		\vspace{0.1cm}
		\textbf{F1 SCORE:} Captures both in a single metric
		
		\textcolor{myred}{\textbf{EXAM TRAP:}} Precision and Recall have \textcolor{myblue}{tradeoff} - cannot maximize both simultaneously!
	\end{frame}
	
	\begin{frame}{T41: F1 Score Definition}
		\fontsize{6.5}{7.5}\selectfont
		\textbf{TOPIC: F1 Score - Balancing Precision and Recall}
		
		\vspace{0.1cm}
		\textbf{FORMULA:}
		\[
		F1 = 2 \times \frac{Precision \times Recall}{Precision + Recall}
		\]
		
		\vspace{0.1cm}
		\textbf{WHY HARMONIC MEAN:}
		\begin{itemize}
			\item Harmonic mean penalizes imbalance
			\item Low precision OR low recall → Low F1
			\item Better than arithmetic mean for this purpose
			\item If either precision or recall is low, F1 is low
		\end{itemize}
		
		\vspace{0.1cm}
		\textbf{INTERPRETATION:}
		\begin{itemize}
			\item Bounded between 0 and 1
			\item Close to 1 = high precision and high recall
			\item Low score = poor precision, poor recall, or both
		\end{itemize}
		
		\vspace{0.1cm}
		\textbf{EXAMPLE (Dividend Prediction):}
		\begin{itemize}
			\item Precision = 78.1\%
			\item Recall = 72.0\%
			\item F1 = 2 × (0.781 × 0.720) / (0.781 + 0.720)
			\item F1 = 2 × 0.5623 / 1.501 = 1.1246 / 1.501 = 74.9\%
		\end{itemize}
		
		\textcolor{myred}{\textbf{EXAM TRAP:}} F1 = \textcolor{myblue}{harmonic mean of precision and recall}!
	\end{frame}
	
	\begin{frame}{T42: F1 Score Example}
		\fontsize{6.5}{7.5}\selectfont
		\textbf{TOPIC: Calculating F1 Score from Precision and Recall}
		
		\vspace{0.1cm}
		\textbf{EXAMPLE DATA:}
		\begin{itemize}
			\item Precision = 78.1\% = 0.781
			\item Recall = 72.0\% = 0.720
		\end{itemize}
		
		\vspace{0.1cm}
		\textbf{CALCULATION:}
		\begin{itemize}
			\item F1 = 2 × (Precision × Recall) / (Precision + Recall)
			\item F1 = 2 × (0.781 × 0.720) / (0.781 + 0.720)
			\item F1 = 2 × 0.5623 / 1.501
			\item F1 = 1.1246 / 1.501
			\item F1 = 0.7494 = 74.9\%
		\end{itemize}
		
		\vspace{0.1cm}
		\textbf{INTERPRETATION:}
		\begin{itemize}
			\item F1 = 74.9\%
			\item Balances precision and recall
			\item Good overall performance
		\end{itemize}
		
		\textcolor{myred}{\textbf{EXAM TRAP:}} F1 = \textcolor{myblue}{2 × (P × R) / (P + R)}!
	\end{frame}
	
	\begin{frame}{T43: Specificity Definition}
		\fontsize{6.5}{7.5}\selectfont
		\textbf{TOPIC: Specificity - True Negative Rate}
		
		\vspace{0.1cm}
		\textbf{FORMULA:}
		\[
		Specificity = \frac{TN}{TN + FP}
		\]
		
		\vspace{0.1cm}
		\textbf{INTERPRETATION:}
		\begin{itemize}
			\item "Of all actual negatives, what percentage did we correctly identify?"
			\item True Negative Rate
			\item Proportion of all negative outcomes correctly predicted as negative
		\end{itemize}
		
		\vspace{0.1cm}
		\textbf{RELATION TO RECALL:}
		\begin{itemize}
			\item Recall = TPR (True Positive Rate)
			\item Specificity = TNR (True Negative Rate)
			\item Both are important for ROC curve
			\item FPR = 1 - Specificity
		\end{itemize}
		
		\vspace{0.1cm}
		\textbf{WHEN SPECIFICITY MATTERS:}
		\begin{itemize}
			\item When false positives are costly
			\item Customer friction concerns
			\item Resource constraints for investigation
		\end{itemize}
		
		\textcolor{myred}{\textbf{EXAM TRAP:}} Specificity = \textcolor{myblue}{TN / (TN + FP)} - True Negative Rate!
	\end{frame}
	
	\begin{frame}{T44: Imbalanced Data Problem}
		\fontsize{6.5}{7.5}\selectfont
		\textbf{TOPIC: Why Accuracy is Misleading for Imbalanced Datasets}
		
		\vspace{0.1cm}
		\textbf{THE PROBLEM:}
		\begin{itemize}
			\item Dataset: 98\% solvent, 2\% default
			\item Model predicts "solvent" for everyone
			\item Accuracy = 98\% (seems great!)
			\item But model catches zero defaults
		\end{itemize}
		
		\vspace{0.1cm}
		\textbf{WHAT METRICS TO USE INSTEAD:}
		\begin{itemize}
			\item \textcolor{myblue}{Precision:} 0\% (no defaults predicted)
			\item \textcolor{myblue}{Recall:} 0\% (caught zero actual defaults)
			\item \textcolor{myblue}{F1:} 0\% (poor precision and recall)
			\item \textcolor{myblue}{AUC:} 0.5 (no discriminatory ability)
		\end{itemize}
		
		\vspace{0.1cm}
		\textbf{EXAMPLE:}
		\begin{itemize}
			\item Fraud detection: 99\% legitimate
			\item Model predicting "legitimate" always has 99\% accuracy
			\item But catches zero fraud (Recall = 0)
		\end{itemize}
		
		\vspace{0.1cm}
		\textbf{LESSON:}
		\begin{itemize}
			\item Accuracy is insufficient for imbalanced data
			\item Use Precision, Recall, F1, AUC
			\item Consider cost-sensitive metrics
		\end{itemize}
		
		\textcolor{myred}{\textbf{EXAM TRAP:}} Accuracy is \textcolor{myblue}{misleading for imbalanced data}!
	\end{frame}
	
	\begin{frame}{T45: Confusion Matrix Example}
		\fontsize{6.5}{7.5}\selectfont
		\textbf{TOPIC: Complete Confusion Matrix Analysis}
		
		\vspace{0.1cm}
		\textbf{EXAMPLE DATA:}
		\begin{itemize}
			\item True Positive (TP) = 60
			\item True Negative (TN) = 498
			\item False Positive (FP) = 98
			\item False Negative (FN) = 5
		\end{itemize}
		
		\vspace{0.1cm}
		\textbf{CALCULATIONS:}
		\begin{itemize}
			\item Total = 60 + 498 + 98 + 5 = 661
			\item Accuracy = (60 + 498) / 661 = 558/661 = 84.4\%
			\item Precision = 60 / (60 + 98) = 60/158 = 38.0\%
			\item Recall = 60 / (60 + 5) = 60/65 = 92.3\%
			\item Error Rate = 1 - 0.844 = 15.6\%
		\end{itemize}
		
		\vspace{0.1cm}
		\textbf{INTERPRETATION:}
		\begin{itemize}
			\item Model catches 92.3\% of defaults (great!)
			\item But only 38\% of default predictions are correct
			\item Many false alarms (98 false positives)
		\end{itemize}
		
		\textcolor{myred}{\textbf{EXAM TRAP:}} Work through each formula carefully - don't mix up denominators!
	\end{frame}
	
	% ============================================================
	% SECTION 5: ROC AND AUC
	% ============================================================
	
	\begin{frame}{T46: ROC Curve Definition}
		\fontsize{6.5}{7.5}\selectfont
		\textbf{TOPIC: Receiver Operating Characteristic (ROC) Curve}
		
		\vspace{0.1cm}
		\textbf{DEFINITION:}
		\begin{itemize}
			\item Graphical illustration of tradeoff between true and false positive rates against different choices of threshold Z
			\item Popular way to visualize classification model performance across all thresholds
		\end{itemize}
		
		\vspace{0.1cm}
		\textbf{ROC CURVE COMPONENTS:}
		\begin{itemize}
			\item \textcolor{myblue}{Y-axis:} True Positive Rate (Recall)
			\item \textcolor{myblue}{X-axis:} False Positive Rate (1 - Specificity)
			\item \textcolor{myblue}{Points:} From varying decision threshold
		\end{itemize}
		
		\vspace{0.1cm}
		\textbf{THRESHOLD EXTREMES:}
		\begin{itemize}
			\item \textcolor{myblue}{Threshold = 0:} All predicted positive
			\begin{itemize}
				\item TPR = 1 (all positives caught)
				\item FPR = 1 (all negatives misclassified)
				\item Point: (1, 1)
			\end{itemize}
			\item \textcolor{myblue}{Threshold = 1:} All predicted negative
			\begin{itemize}
				\item TPR = 0 (no positives caught)
				\item FPR = 0 (all negatives correct)
				\item Point: (0, 0)
			\end{itemize}
		\end{itemize}
		
		\vspace{0.1cm}
		\textbf{KEY INSIGHT:}
		\begin{itemize}
			\item For values between 0 and 1, TPR increases when FPR increases
			\item Shows tradeoff between catching positives and avoiding false alarms
		\end{itemize}
		
		\textcolor{myred}{\textbf{EXAM TRAP:}} ROC shows \textcolor{myblue}{TPR vs FPR} tradeoff across thresholds!
	\end{frame}
	
	\begin{frame}{T47: AUC Interpretation}
		\fontsize{6.5}{7.5}\selectfont
		\textbf{TOPIC: Area Under ROC Curve (AUC) - Interpretation}
		
		\vspace{0.1cm}
		\textbf{DEFINITION:}
		\begin{itemize}
			\item Area under the ROC curve
			\item Measures overall ability of model to distinguish between positive and negative classes
			\item Threshold-independent measure
		\end{itemize}
		
		\vspace{0.1cm}
		\textbf{AUC VALUES AND INTERPRETATION:}
		\begin{itemize}
			\item \textcolor{mygreen}{AUC = 1.0:} Perfect discrimination
			\item \textcolor{myblue}{AUC = 0.5:} No discrimination (random guessing)
			\item \textcolor{myred}{AUC < 0.5:} Worse than random (should invert predictions)
			\item \textcolor{mygreen}{Higher AUC = better model}
		\end{itemize}
		
		\vspace{0.1cm}
		\textbf{PROBABILISTIC INTERPRETATION:}
		\begin{itemize}
			\item Probability that a random positive ranks higher than a random negative
			\item Measures ranking ability of model
			\item Threshold-independent
		\end{itemize}
		
		\vspace{0.1cm}
		\textbf{EXAMPLE:}
		\begin{itemize}
			\item AUC = 0.8 means 80\% chance that a randomly chosen positive is ranked higher than a randomly chosen negative
		\end{itemize}
		
		\textcolor{myred}{\textbf{EXAM TRAP:}} AUC = \textcolor{myblue}{probability that positive ranks higher than negative}!
	\end{frame}
	
	\begin{frame}{T48: AUC = 0.5 Meaning}
		\fontsize{6.5}{7.5}\selectfont
		\textbf{TOPIC: What AUC = 0.5 Indicates}
		
		\vspace{0.1cm}
		\textbf{AUC = 0.5:}
		\begin{itemize}
			\item No predictive ability
			\item Model cannot distinguish between classes
			\item Equivalent to random guessing
			\item Dashed diagonal line on ROC plot
		\end{itemize}
		
		\vspace{0.1cm}
		\textbf{WHAT IT MEANS:}
		\begin{itemize}
			\item Model has no useful predictive information
			\item Positive and negative cases are equally likely to be ranked higher
			\item No better than flipping a coin
		\end{itemize}
		
		\vspace{0.1cm}
		\textbf{AUC < 0.5:}
		\begin{itemize}
			\item Worse than random predictions
			\item Model performs worse than random draws
			\item Should invert predictions (swap positive/negative)
		\end{itemize}
		
		\vspace{0.1cm}
		\textbf{TYPICAL GOOD MODELS:}
		\begin{itemize}
			\item AUC between 0.7 and 0.9
			\item AUC > 0.9: Excellent discrimination
		\end{itemize}
		
		\textcolor{myred}{\textbf{EXAM TRAP:}} AUC = 0.5 = \textcolor{myblue}{no predictive ability} - random guessing!
	\end{frame}
	
	\begin{frame}{T49: AUC = 1.0 Meaning}
		\fontsize{6.5}{7.5}\selectfont
		\textbf{TOPIC: What AUC = 1.0 Indicates}
		
		\vspace{0.1cm}
		\textbf{AUC = 1.0:}
		\begin{itemize}
			\item Perfect predictive ability
			\item Perfect separation between classes
			\item Every positive has higher score than every negative
			\item There exists a threshold with perfect classification
		\end{itemize}
		
		\vspace{0.1cm}
		\textbf{CAUTION:}
		\begin{itemize}
			\item AUC = 1.0 on test data is suspicious
			\item Likely indicates overfitting or data leakage
			\item Real-world data rarely has AUC = 1.0
			\item Should investigate immediately
		\end{itemize}
		
		\vspace{0.1cm}
		\textbf{WHAT TO CHECK:}
		\begin{itemize}
			\item Is there data leakage from test to training?
			\item Is the model too complex?
			\item Are there duplicate records?
			\item Is there temporal leakage?
		\end{itemize}
		
		\textcolor{myred}{\textbf{EXAM TRAP:}} AUC = 1.0 = \textcolor{myblue}{perfect separation} - be suspicious!
	\end{frame}
	
	\begin{frame}{T50: ROC Curve Shape}
		\fontsize{6.5}{7.5}\selectfont
		\textbf{TOPIC: Interpreting ROC Curve Shape}
		
		\vspace{0.1cm}
		\textbf{BETTER ROC CURVE:}
		\begin{itemize}
			\item Curve closer to upper-left corner
			\item Higher TPR, lower FPR
			\item More area under curve (higher AUC)
		\end{itemize}
		
		\vspace{0.1cm}
		\textbf{ROC CURVE INTERPRETATION:}
		\begin{itemize}
			\item \textcolor{mygreen}{Upper-left corner:} TPR=1, FPR=0 (perfect)
			\item \textcolor{myblue}{Diagonal line:} TPR=FPR (random)
			\item \textcolor{myred}{Lower-right corner:} TPR=0, FPR=1 (worst)
		\end{itemize}
		
		\vspace{0.1cm}
		\textbf{COMPARISON:}
		\begin{itemize}
			\item Curve A with more area = better
			\item Curve that stays above and left = better
			\item AUC captures this visually
		\end{itemize}
		
		\vspace{0.1cm}
		\textbf{SAMPLE ROC CURVE:}
		\begin{itemize}
			\item Model A: Curve closer to upper-left (better)
			\item Model B: Curve closer to diagonal (worse)
			\item Model A has higher AUC
		\end{itemize}
		
		\textcolor{myred}{\textbf{EXAM TRAP:}} Better ROC curve = \textcolor{myblue}{closer to upper-left corner}!
	\end{frame}
	
	\begin{frame}{T51: Threshold Effect on ROC}
		\fontsize{6.5}{7.5}\selectfont
		\textbf{TOPIC: How Threshold Changes Affect ROC Points}
		
		\vspace{0.1cm}
		\textbf{THRESHOLD EXTREMES:}
		\begin{itemize}
			\item \textcolor{myblue}{Threshold = 0:} All predicted positive
			\begin{itemize}
				\item TPR = 1 (all positives caught)
				\item FPR = 1 (all negatives misclassified)
				\item Point: (1, 1)
			\end{itemize}
			\item \textcolor{myblue}{Threshold = 1:} All predicted negative
			\begin{itemize}
				\item TPR = 0 (no positives caught)
				\item FPR = 0 (all negatives correct)
				\item Point: (0, 0)
			\end{itemize}
		\end{itemize}
		
		\vspace{0.1cm}
		\textbf{AS THRESHOLD INCREASES:}
		\begin{itemize}
			\item Points move from (1,1) toward (0,0)
			\item TPR decreases (catch fewer positives)
			\item FPR decreases (fewer false alarms)
		\end{itemize}
		
		\vspace{0.1cm}
		\textbf{AS THRESHOLD DECREASES:}
		\begin{itemize}
			\item Points move from (0,0) toward (1,1)
			\item TPR increases (catch more positives)
			\item FPR increases (more false alarms)
		\end{itemize}
		
		\textcolor{myred}{\textbf{EXAM TRAP:}} ROC goes from \textcolor{myblue}{(0,0) to (1,1)} as threshold decreases!
	\end{frame}
	
	\begin{frame}{T52: AUC for Model Comparison}
		\fontsize{6.5}{7.5}\selectfont
		\textbf{TOPIC: Using AUC to Compare Models in Lending Decisions}
		
		\vspace{0.1cm}
		\textbf{AUC IN LENDING:}
		\begin{itemize}
			\item Measures ability to rank borrowers by risk
			\item Higher AUC = better separation of good from bad loans
			\item Better model = more accurate risk ranking
		\end{itemize}
		
		\vspace{0.1cm}
		\textbf{DECISION CONTEXT:}
		\begin{itemize}
			\item Choose model with higher AUC on test set
			\item AUC = 0.8 vs 0.7: First has better ranking
			\item Better AUC = more confidence in decisions
		\end{itemize}
		
		\vspace{0.1cm}
		\textbf{HOW IT WORKS:}
		\begin{itemize}
			\item Model ranks all loan applicants by default probability
			\item Higher AUC means better able to separate good from bad
			\item Can set threshold based on risk appetite
		\end{itemize}
		
		\vspace{0.1cm}
		\textbf{EXAMPLE:}
		\begin{itemize}
			\item Model A: AUC = 0.85
			\item Model B: AUC = 0.72
			\item Model A better at ranking borrowers by risk
			\item Choose Model A for lending decisions
		\end{itemize}
		
		\textcolor{myred}{\textbf{EXAM TRAP:}} Higher AUC = \textcolor{myblue}{better model} for lending decisions!
	\end{frame}
	
	\begin{frame}{T53: Precision-Recall vs ROC}
		\fontsize{6.5}{7.5}\selectfont
		\textbf{TOPIC: When to Use Precision-Recall Curve Instead of ROC}
		
		\vspace{0.1cm}
		\textbf{PRECISION-RECALL CURVES:}
		\begin{itemize}
			\item Better for imbalanced datasets
			\item Focus on positive class performance
			\item More sensitive to rare events
		\end{itemize}
		
		\vspace{0.1cm}
		\textbf{WHY NOT ROC:}
		\begin{itemize}
			\item ROC can be optimistic with imbalance
			\item Large number of negatives makes FPR low
			\item AUC can be high even with poor positive performance
		\end{itemize}
		
		\vspace{0.1cm}
		\textbf{WHEN TO USE EACH:}
		\begin{itemize}
			\item \textcolor{myblue}{ROC:} Balanced datasets or when both classes equally important
			\item \textcolor{myblue}{Precision-Recall:} Imbalanced datasets, focus on positive class
		\end{itemize}
		
		\vspace{0.1cm}
		\textbf{EXAMPLE:}
		\begin{itemize}
			\item Fraud detection (1\% fraud): Use Precision-Recall
			\item Credit approval (50/50): Use ROC
		\end{itemize}
		
		\textcolor{myred}{\textbf{EXAM TRAP:}} Precision-Recall for \textcolor{myblue}{imbalanced data}!
	\end{frame}
	
	\begin{frame}{T54: Multi-class ROC}
		\fontsize{6.5}{7.5}\selectfont
		\textbf{TOPIC: Extending Classification Metrics to Multi-class Problems}
		
		\vspace{0.1cm}
		\textbf{ONE-VS-REST APPROACH:}
		\begin{itemize}
			\item For each class, treat as "positive"
			\item All other classes as "negative"
			\item Compute TP, FP, FN, TN for each class
			\item Average or aggregate metrics
		\end{itemize}
		
		\vspace{0.1cm}
		\textbf{EXAMPLE - CREDIT RATINGS:}
		\begin{itemize}
			\item Classes: AAA, AA, A, BBB, BB, B, CCC
			\item For AAA: positive = AAA, negative = all others
			\item Compute precision/recall for AAA
			\item Repeat for each rating
		\end{itemize}
		
		\vspace{0.1cm}
		\textbf{GENERALIZATION:}
		\begin{itemize}
			\item Precision: Sum all TP / Sum all TP + Sum all FP
			\item Recall: Sum all TP / Sum all TP + Sum all FN
			\item Other metrics generalized similarly
		\end{itemize}
		
		\vspace{0.1cm}
		\textbf{APPLICATION:}
		\begin{itemize}
			\item Multi-class credit rating prediction
			\item Risk categorization (high/medium/low)
			\item Severity classification
		\end{itemize}
		
		\textcolor{myred}{\textbf{EXAM TRAP:}} Multi-class uses \textcolor{myblue}{one-vs-rest} approach!
	\end{frame}
	
	\begin{frame}{T55: AUC vs Accuracy}
		\fontsize{6.5}{7.5}\selectfont
		\textbf{TOPIC: Advantages of AUC Over Accuracy}
		
		\vspace{0.1cm}
		\textbf{AUC ADVANTAGES:}
		\begin{itemize}
			\item \textcolor{myblue}{Threshold-independent:} Evaluates all thresholds
			\item \textcolor{myblue}{Works with imbalance:} Not skewed by majority class
			\item \textcolor{myblue}{Ranking ability:} Measures how well model orders cases
			\item \textcolor{myblue}{Comprehensive:} Single number captures overall performance
		\end{itemize}
		
		\vspace{0.1cm}
		\textbf{ACCURACY LIMITATIONS:}
		\begin{itemize}
			\item Threshold-dependent (uses one threshold)
			\item Misleading for imbalanced data
			\item Doesn't capture ranking ability
			\item Treats all errors equally
		\end{itemize}
		
		\vspace{0.1cm}
		\textbf{EXAMPLE:}
		\begin{itemize}
			\item Dataset: 98\% solvent, 2\% default
			\item Model predicting all solvent: 98\% accuracy
			\item But AUC = 0.5 (no discrimination)
			\item Accuracy misleading, AUC reveals truth
		\end{itemize}
		
		\textcolor{myred}{\textbf{EXAM TRAP:}} AUC is \textcolor{myblue}{threshold-independent} - a key advantage!
	\end{frame}
	
	% ============================================================
	% SECTION 6: THRESHOLD SELECTION
	% ============================================================
	
	\begin{frame}{T56: Threshold Definition}
		\fontsize{6.5}{7.5}\selectfont
		\textbf{TOPIC: Decision Threshold in Classification}
		
		\vspace{0.1cm}
		\textbf{DEFINITION:}
		\begin{itemize}
			\item Most classification models output a probability or score
			\item A decision threshold is applied to convert this score into a binary classification
			\item Example: If PD > 5\%, classify as "high risk"; if fraud score > 0.7, block transaction
		\end{itemize}
		
		\vspace{0.1cm}
		\textbf{COMMON THRESHOLDS:}
		\begin{itemize}
			\item \textcolor{myblue}{Default:} 0.5 (balanced)
			\item \textcolor{myblue}{Can be adjusted:} Based on business needs
			\item \textcolor{myblue}{Different thresholds:} Different tradeoffs
		\end{itemize}
		
		\vspace{0.1cm}
		\textbf{WHY THRESHOLD CHOICE MATTERS:}
		\begin{itemize}
			\item Directly determines balance of False Positives and False Negatives
			\item For every threshold, different Confusion Matrix
			\item Different Precision and Recall values
			\item Critical business decision, not just technical
		\end{itemize}
		
		\vspace{0.1cm}
		\textbf{EXAMPLE:}
		\begin{itemize}
			\item Model outputs probability of default
			\item Threshold = 0.5: Predict default if PD > 0.5
			\item Threshold = 0.3: More conservative, predict more defaults
		\end{itemize}
		
		\textcolor{myred}{\textbf{EXAM TRAP:}} Threshold converts \textcolor{myblue}{probability to class label}!
	\end{frame}
	
	\begin{frame}{T57: Cost Matrix Definition}
		\fontsize{6.5}{7.5}\selectfont
		\textbf{TOPIC: Cost Matrix for Threshold Selection}
		
		\vspace{0.1cm}
		\textbf{DEFINITION:}
		\begin{itemize}
			\item Explicitly defining economic consequences associated with each quadrant of Confusion Matrix
			\item Most effective way to guide threshold selection
		\end{itemize}
		
		\vspace{0.1cm}
		\textbf{COST MATRIX COMPONENTS:}
		\begin{itemize}
			\item \textcolor{myblue}{Cost of False Negative:} Average loss from undetected fraud; capital impact of missed default (often largest cost)
			\item \textcolor{myblue}{Cost of False Positive:} Operational cost to investigate false alert; potential lost revenue from declining good customer
			\item \textcolor{myblue}{Benefit of True Positive:} Loss avoided by catching fraud; benefit of early intervention
			\item \textcolor{myblue}{Benefit of True Negative:} Smooth processing; benefit of correctly underwriting good loan
		\end{itemize}
		
		\vspace{0.1cm}
		\textbf{OPTIMAL THRESHOLD:}
		\begin{itemize}
			\item Threshold that minimizes total expected cost
			\item OR maximizes total expected benefit
			\item Once costs/benefits estimated (even roughly)
		\end{itemize}
		
		\textcolor{myred}{\textbf{EXAM TRAP:}} Cost matrix = \textcolor{myblue}{economic consequences} of each decision!
	\end{frame}
	
	\begin{frame}{T58: Stakeholder Input}
		\fontsize{6.5}{7.5}\selectfont
		\textbf{TOPIC: Cross-Functional Input for Cost Matrix}
		
		\vspace{0.1cm}
		\textbf{KEY STAKEHOLDERS:}
		\begin{itemize}
			\item \textcolor{myblue}{Business Lines:} Understand cost of missed opportunities, customer attrition
			\item \textcolor{myblue}{Operations:} Know tangible costs of manual reviews (investigator time, call center)
			\item \textcolor{myblue}{Finance/Risk:} Can quantify financial impact of missed risks
			\item \textcolor{myblue}{Compliance/Legal:} Advise on potential costs of regulatory breaches
		\end{itemize}
		
		\vspace{0.1cm}
		\textbf{WHY MULTIPLE PERSPECTIVES:}
		\begin{itemize}
			\item Different stakeholders understand different costs
			\item Complete picture of all consequences
			\item Collaborative decisions are more defensible
			\item Developing cost matrix is not purely quantitative exercise
		\end{itemize}
		
		\vspace{0.1cm}
		\textbf{EXAMPLES:}
		\begin{itemize}
			\item Business: Cost of declining good loan
			\item Operations: Cost of AML alert investigation
			\item Finance: Capital impact of missed default
			\item Legal: Cost of regulatory breach
		\end{itemize}
		
		\textcolor{myred}{\textbf{EXAM TRAP:}} Cost matrix needs \textcolor{myblue}{cross-functional input}!
	\end{frame}
	
	\begin{frame}{T59: Dynamic Threshold Adjustment}
		\fontsize{6.5}{7.5}\selectfont
		\textbf{TOPIC: Thresholds Should Be Dynamic, Not Static}
		
		\vspace{0.1cm}
		\textbf{DYNAMIC ADJUSTMENT EXAMPLES:}
		\begin{itemize}
			\item \textcolor{myblue}{Holiday Fraud Spike:} Lower threshold (higher recall, catch more fraud)
			\item \textcolor{myblue}{Economic Downturn:} Adjust default thresholds
			\item \textcolor{myblue}{Benign Periods:} Higher threshold (fewer false positives)
		\end{itemize}
		
		\vspace{0.1cm}
		\textbf{PREREQUISITES:}
		\begin{itemize}
			\item Clear risk appetite statement
			\item Documented review process
			\item Transparent justification
			\item Regulatory awareness
		\end{itemize}
		
		\vspace{0.1cm}
		\textbf{DOCUMENTATION EXAMPLE:}
		\begin{itemize}
			\item "Threshold set for AML monitoring to achieve 80\% recall"
			\item "Acceptable false positive rate of X\%"
			\item "Deemed acceptable given consequences of missing AML events"
		\end{itemize}
		
		\vspace{0.1cm}
		\textbf{CAUTION:}
		\begin{itemize}
			\item Ensure threshold adjusted for right reasons
			\item Not "set and forget"
			\item Actively managed
		\end{itemize}
		
		\textcolor{myred}{\textbf{EXAM TRAP:}} Thresholds should be \textcolor{myblue}{dynamic} - not "set and forget"!
	\end{frame}
	
	\begin{frame}{T60: Threshold Documentation}
		\fontsize{6.5}{7.5}\selectfont
		\textbf{TOPIC: Documentation and Defensibility of Threshold Choices}
		
		\vspace{0.1cm}
		\textbf{DOCUMENTATION SHOULD SHOW:}
		\begin{itemize}
			\item How threshold aligns with risk appetite
			\item Tradeoff between recall and false positives
			\item Costs and benefits considered
			\item Periodic review process
		\end{itemize}
		
		\vspace{0.1cm}
		\textbf{WHY IT MATTERS:}
		\begin{itemize}
			\item Regulatory requirement
			\item Audit trail
			\item Knowledge transfer
			\item Defensible decisions
		\end{itemize}
		
		\vspace{0.1cm}
		\textbf{EXAMPLE STATEMENT:}
		\begin{itemize}
			\item "We are currently setting threshold for transaction monitoring system to achieve 80\% recall for high-risk AML typologies"
			\item "This means we accept a False Positive Rate of X\%, costing Y in operational overhead"
			\item "Deemed acceptable given severe consequences of missing AML events"
		\end{itemize}
		
		\vspace{0.1cm}
		\textbf{PERIODIC REVIEW:}
		\begin{itemize}
			\item Demonstrates not "set and forget"
			\item Actively managed in response to changing conditions
			\item Documented review process
		\end{itemize}
		
		\textcolor{myred}{\textbf{EXAM TRAP:}} Thresholds must be \textcolor{myblue}{transparent and defensible}!
	\end{frame}
	
	\begin{frame}{T61: False Positive Cost}
		\fontsize{6.5}{7.5}\selectfont
		\textbf{TOPIC: Costs Associated with False Positives in AML Monitoring}
		
		\vspace{0.1cm}
		\textbf{FALSE POSITIVE COSTS:}
		\begin{enumerate}
			\item \textcolor{myblue}{Operational Cost:}
			\begin{itemize}
				\item Investigator time wasted
				\item Compliance teams chasing red herrings
				\item Resource consumption
			\end{itemize}
			
			\item \textcolor{myblue}{Customer Friction:}
			\begin{itemize}
				\item Declined legitimate transactions
				\item Customer inconvenience
				\item Potential customer attrition
			\end{itemize}
			
			\item \textcolor{myblue}{Reputational Cost:}
			\begin{itemize}
				\item Customer dissatisfaction
				\item Negative word-of-mouth
			\end{itemize}
			
			\item \textcolor{myblue}{Alert Fatigue:}
			\begin{itemize}
				\item Teams desensitized to alerts
				\item May miss real issues
			\end{itemize}
		\end{enumerate}
		
		\vspace{0.1cm}
		\textbf{WHY THIS MATTERS:}
		\begin{itemize}
			\item High FP rate = inefficient operations
			\item Can outweigh benefits of catching fraud
			\item Affects customer experience
		\end{itemize}
		
		\textcolor{myred}{\textbf{EXAM TRAP:}} False Positives = \textcolor{myblue}{operational and customer costs}!
	\end{frame}
	
	\begin{frame}{T62: False Negative Cost}
		\fontsize{6.5}{7.5}\selectfont
		\textbf{TOPIC: Costs Associated with False Negatives in Default Prediction}
		
		\vspace{0.1cm}
		\textbf{FALSE NEGATIVE COSTS:}
		\begin{enumerate}
			\item \textcolor{myblue}{Financial Losses:}
			\begin{itemize}
				\item Actual loan losses
				\item Direct financial impact
				\item Hard to recover
			\end{itemize}
			
			\item \textcolor{myblue}{Regulatory Capital Implications:}
			\begin{itemize}
				\item Capital add-ons
				\item Increased capital requirements
				\item Regulatory scrutiny
			\end{itemize}
			
			\item \textcolor{myblue}{Risk Management:}
			\begin{itemize}
				\item Misclassification of risk
				\item Underpricing of risk
			\end{itemize}
			
			\item \textcolor{myblue}{Reputational Impact:}
			\begin{itemize}
				\item Regulatory sanctions
				\item Market confidence erosion
			\end{itemize}
		\end{enumerate}
		
		\vspace{0.1cm}
		\textbf{WHY OFTEN LARGEST COST:}
		\begin{itemize}
			\item Direct financial impact
			\item Hard to recover from
			\item Regulatory consequences
		\end{itemize}
		
		\textcolor{myred}{\textbf{EXAM TRAP:}} False Negatives = \textcolor{myblue}{financial and regulatory costs} - usually largest!
	\end{frame}
	
	\begin{frame}{T63: Optimal Threshold Selection}
		\fontsize{6.5}{7.5}\selectfont
		\textbf{TOPIC: Determining the Optimal Threshold Using Cost Matrix}
		
		\vspace{0.1cm}
		\textbf{EXPECTED COST CALCULATION:}
		\begin{itemize}
			\item For each threshold, compute expected cost
			\item Cost = C(TP)×P(TP) + C(FP)×P(FP) + C(FN)×P(FN) + C(TN)×P(TN)
			\item Choose threshold with minimum cost
			\item Or maximum expected benefit
		\end{itemize}
		
		\vspace{0.1cm}
		\textbf{COST-BASED THRESHOLD:}
		\begin{itemize}
			\item Different from accuracy-based threshold
			\item Reflects business reality
			\item More defensible
		\end{itemize}
		
		\vspace{0.1cm}
		\textbf{PROCESS:}
		\begin{enumerate}
			\item Define costs/benefits for each quadrant
			\item For each possible threshold, compute expected cost
			\item Select threshold minimizing cost
			\item Document decision
		\end{enumerate}
		
		\textcolor{myred}{\textbf{EXAM TRAP:}} Optimal threshold = \textcolor{myblue}{minimum expected cost}!
	\end{frame}
	
	\begin{frame}{T64: Threshold and Business Needs}
		\fontsize{6.5}{7.5}\selectfont
		\textbf{TOPIC: Which Model for Fraud Detection?}
		
		\vspace{0.1cm}
		\textbf{FRAUD DETECTION PRIORITIES:}
		\begin{itemize}
			\item Missing fraud (False Negative) is very costly
			\item Better to catch more fraud (high recall)
			\item Accept some false alarms (lower precision)
			\item Business needs drive the tradeoff
		\end{itemize}
		
		\vspace{0.1cm}
		\textbf{OTHER CONTEXTS:}
		\begin{itemize}
			\item \textcolor{myblue}{AML Monitoring:} High recall critical
			\item \textcolor{myblue}{Credit Approval:} Balance needed
			\item \textcolor{myblue}{Marketing:} Precision more important
		\end{itemize}
		
		\vspace{0.1cm}
		\textbf{EXAMPLE:}
		\begin{itemize}
			\item Model A: High precision (90\%), low recall (50\%)
			\item Model B: Lower precision (60\%), high recall (90\%)
			\item For fraud detection, Model B preferred
			\item Catches more fraud despite more false alarms
		\end{itemize}
		
		\textcolor{myred}{\textbf{EXAM TRAP:}} Fraud detection prioritizes \textcolor{myblue}{recall over precision}!
	\end{frame}
	
	\begin{frame}{T65: Threshold Review Process}
		\fontsize{6.5}{7.5}\selectfont
		\textbf{TOPIC: How Often Should Thresholds Be Reviewed?}
		
		\vspace{0.1cm}
		\textbf{REVIEW TRIGGERS:}
		\begin{itemize}
			\item Significant environmental changes
			\item Shifts in risk appetite
			\item New regulatory requirements
			\item Changes in fraud patterns
			\item Economic downturns
		\end{itemize}
		
		\vspace{0.1cm}
		\textbf{BEST PRACTICE:}
		\begin{itemize}
			\item Documented review schedule
			\item Clear triggers for review
			\item Stakeholder involvement
			\item Regulatory notification when appropriate
		\end{itemize}
		
		\vspace{0.1cm}
		\textbf{FREQUENCY:}
		\begin{itemize}
			\item \textcolor{myblue}{Regular:} Quarterly, semi-annually, annually
			\item \textcolor{myblue}{Event-driven:} As needed based on triggers
			\item \textcolor{myblue}{Not:} Never or "set and forget"
		\end{itemize}
		
		\textcolor{myred}{\textbf{EXAM TRAP:}} Thresholds should be \textcolor{myblue}{periodically reviewed} - not static!
	\end{frame}
	
	% ============================================================
	% SECTION 7: MODEL COMPARISON
	% ============================================================
	
	\begin{frame}{T66: Model Comparison Example}
		\fontsize{6.5}{7.5}\selectfont
		\textbf{TOPIC: Comparing Logistic Regression and Neural Network}
		
		\vspace{0.1cm}
		\textbf{EXAMPLE: LendingClub Default Prediction}
		
		\begin{tabular}{|p{2.5cm}|p{2.5cm}|p{2.5cm}|p{2.5cm}|p{2.5cm}|}
			\hline
			\textbf{Measure} & \textbf{LR Train} & \textbf{NN Train} & \textbf{LR Test} & \textbf{NN Test} \\
			\hline
			Accuracy & 0.842 & 0.842 & 0.713 & 0.701 \\
			\hline
			Precision & 0.656 & 0.750 & 0.600 & 0.541 \\
			\hline
			Recall & 0.236 & 0.169 & 0.283 & 0.377 \\
			\hline
		\end{tabular}
		
		\vspace{0.1cm}
		\textbf{KEY OBSERVATIONS:}
		\begin{itemize}
			\item Performance drops from training to test (some overfitting)
			\item Different metrics give different rankings
			\item \textcolor{myblue}{Logistic Regression:} Better accuracy and precision on test
			\item \textcolor{myblue}{Neural Network:} Better recall on test
			\item No clear winner; choice depends on priorities
		\end{itemize}
		
		\textcolor{myred}{\textbf{EXAM TRAP:}} Different metrics = \textcolor{myblue}{different model rankings}!
	\end{frame}
	
	\begin{frame}{T67: Training vs Test Comparison}
		\fontsize{6.5}{7.5}\selectfont
		\textbf{TOPIC: What Does Drop from Training to Test Indicate?}
		
		\vspace{0.1cm}
		\textbf{PERFORMANCE DROPS:}
		\begin{itemize}
			\item Training accuracy: 84.2\% → Test: 71.3\% (LR)
			\item Training accuracy: 84.2\% → Test: 70.1\% (NN)
			\item Both show significant decrease
		\end{itemize}
		
		\vspace{0.1cm}
		\textbf{INTERPRETATION:}
		\begin{itemize}
			\item Models overfit to training data
			\item Reasonable amount of overfitting
			\item Could improve with regularization
			\item Remove irrelevant features
		\end{itemize}
		
		\vspace{0.1cm}
		\textbf{REMEDIATION:}
		\begin{itemize}
			\item "Remove some of the least empirically relevant features"
			\item "Apply regularization to the fitted models"
		\end{itemize}
		
		\textcolor{myred}{\textbf{EXAM TRAP:}} Performance drop = \textcolor{myblue}{overfitting} - expected but should be managed!
	\end{frame}
	
	\begin{frame}{T68: Neural Network Overfitting}
		\fontsize{6.5}{7.5}\selectfont
		\textbf{TOPIC: Checking for Overfitting in Neural Networks}
		
		\vspace{0.1cm}
		\textbf{OVERFITTING CHECKS:}
		\begin{enumerate}
			\item \textcolor{myblue}{Performance:} Training >> Test
			\item \textcolor{myblue}{Weights:} Very large or offsetting values
			\item \textcolor{myblue}{Loss:} Training loss flat, validation increasing
		\end{enumerate}
		
		\vspace{0.1cm}
		\textbf{WHAT TO LOOK FOR:}
		\begin{itemize}
			\item Large positive and negative weights
			\item Model memorizing training data
			\item Poor generalization
			\item Strong or weak connections
		\end{itemize}
		
		\vspace{0.1cm}
		\textbf{EXAMPLE:}
		\begin{itemize}
			\item Neural network contains ten units in hidden layer
			\item Uses logistic activation function
			\item Examine fitted weights for offsetting connections
			\item Indicative of overfitting
		\end{itemize}
		
		\textcolor{myred}{\textbf{EXAM TRAP:}} Check \textcolor{myblue}{both performance gap and weight patterns} for overfitting!
	\end{frame}
	
	\begin{frame}{T69: Logistic Regression Advantages}
		\fontsize{6.5}{7.5}\selectfont
		\textbf{TOPIC: Advantages of Logistic Regression Over Neural Networks}
		
		\vspace{0.1cm}
		\textbf{LOGISTIC REGRESSION ADVANTAGES:}
		\begin{itemize}
			\item \textcolor{myblue}{Interpretability:} Coefficients explain feature impact
			\item \textcolor{myblue}{Explainability:} Easy to communicate to stakeholders
			\item \textcolor{myblue}{Speed:} Faster to train and deploy
			\item \textcolor{myblue}{Transparency:} Regulatory acceptance
		\end{itemize}
		
		\vspace{0.1cm}
		\textbf{NEURAL NETWORK ADVANTAGES:}
		\begin{itemize}
			\item Can capture complex patterns
			\item Higher potential accuracy
			\item Better for large datasets
		\end{itemize}
		
		\vspace{0.1cm}
		\textbf{PRACTICAL CONSIDERATIONS:}
		\begin{itemize}
			\item In LendingClub example, LR better on test accuracy and precision
			\item Simpler model may be preferred
			\item Regulatory requirements may favor LR
		\end{itemize}
		
		\textcolor{myred}{\textbf{EXAM TRAP:}} Logistic regression = \textcolor{myblue}{more interpretable} than neural networks!
	\end{frame}
	
	\begin{frame}{T70: Neural Network Advantages}
		\fontsize{6.5}{7.5}\selectfont
		\textbf{TOPIC: Advantages of Neural Networks Over Logistic Regression}
		
		\vspace{0.1cm}
		\textbf{NEURAL NETWORK ADVANTAGES:}
		\begin{itemize}
			\item \textcolor{myblue}{Complex Patterns:} Can model nonlinear relationships
			\item \textcolor{myblue}{Flexibility:} Universal approximator
			\item \textcolor{myblue}{Performance:} Often better for complex tasks
			\item \textcolor{myblue}{Feature Learning:} Learns representations automatically
		\end{itemize}
		
		\vspace{0.1cm}
		\textbf{DISADVANTAGES:}
		\begin{itemize}
			\item Less interpretable
			\item More data needed
			\item More computationally expensive
			\item Risk of overfitting
		\end{itemize}
		
		\vspace{0.1cm}
		\textbf{IN LENDINGCLUB EXAMPLE:}
		\begin{itemize}
			\item Neural network had better recall (37.7\% vs 28.3\%)
			\item Caught more defaults
			\item But lower accuracy and precision
			\item Tradeoff between metrics
		\end{itemize}
		
		\textcolor{myred}{\textbf{EXAM TRAP:}} Neural networks = \textcolor{myblue}{better for complex patterns} but less interpretable!
	\end{frame}
	
	\begin{frame}{T71: Model Selection Criteria}
		\fontsize{6.5}{7.5}\selectfont
		\textbf{TOPIC: In LendingClub, Which Model Would a Risk Manager Choose?}
		
		\vspace{0.1cm}
		\textbf{TRADEOFFS:}
		\begin{itemize}
			\item \textcolor{myblue}{If catching defaults is critical:} Neural network (higher recall)
			\item \textcolor{myblue}{If interpretability is required:} Logistic regression
			\item \textcolor{myblue}{If accuracy is priority:} Logistic regression (slightly higher)
			\item \textcolor{myblue}{If regulatory compliance:} Logistic regression
		\end{itemize}
		
		\vspace{0.1cm}
		\textbf{KEY INSIGHT:}
		\begin{itemize}
			\item No single "best" model
			\item Choice depends on business needs
			\item Must align with risk appetite
			\item Contradictory indicators illustrate importance of choosing appropriate metric
		\end{itemize}
		
		\vspace{0.1cm}
		\textbf{DECISION FRAMEWORK:}
		\begin{itemize}
			\item What is the cost of missing a default?
			\item What is the cost of false alarm?
			\item What are regulatory requirements?
			\item What is stakeholder preference?
		\end{itemize}
		
		\textcolor{myred}{\textbf{EXAM TRAP:}} Model choice = \textcolor{myblue}{depends on priorities and constraints}!
	\end{frame}
	
	\begin{frame}{T72: Overfitting Remediation}
		\fontsize{6.5}{7.5}\selectfont
		\textbf{TOPIC: How to Address Overfitting in LendingClub Example}
		
		\vspace{0.1cm}
		\textbf{REMEDIATION STRATEGIES:}
		\begin{enumerate}
			\item \textcolor{myblue}{Simplify Model:} Remove irrelevant features
			\item \textcolor{myblue}{Add Regularization:} L1, L2, dropout
			\item \textcolor{myblue}{Get More Data:} More training samples
			\item \textcolor{myblue}{Cross-Validation:} Better model selection
			\item \textcolor{myblue}{Early Stopping:} Stop before overfitting
		\end{enumerate}
		
		\vspace{0.1cm}
		\textbf{CHAPTER RECOMMENDATION:}
		\begin{itemize}
			\item "Remove some of the least empirically relevant features"
			\item "Apply regularization to the fitted models"
		\end{itemize}
		
		\vspace{0.1cm}
		\textbf{WHY THIS WORKS:}
		\begin{itemize}
			\item Reduces model complexity
			\item Prevents memorization
			\item Improves generalization
		\end{itemize}
		
		\textcolor{myred}{\textbf{EXAM TRAP:}} Fix overfitting with \textcolor{myblue}{simplification or regularization}!
	\end{frame}
	
	\begin{frame}{T73: Model Comparison Metrics}
		\fontsize{6.5}{7.5}\selectfont
		\textbf{TOPIC: Which Metrics Should Be Used to Compare Models for Default Prediction?}
		
		\vspace{0.1cm}
		\textbf{RECOMMENDED APPROACH:}
		\begin{enumerate}
			\item Start with confusion matrix
			\item Compute Accuracy, Precision, Recall, F1
			\item Consider business costs
			\item Choose metric that aligns with priorities
		\end{enumerate}
		
		\vspace{0.1cm}
		\textbf{WHY MULTIPLE METRICS:}
		\begin{itemize}
			\item Single metric masks important differences
			\item Different stakeholders need different views
			\item Regulatory requirements may specify
			\item No single metric tells whole story
		\end{itemize}
		
		\vspace{0.1cm}
		\textbf{EXAMPLE:}
		\begin{itemize}
			\item LR: Better accuracy and precision
			\item NN: Better recall
			\item Both have merits
			\item Choice depends on what matters most
		\end{itemize}
		
		\textcolor{myred}{\textbf{EXAM TRAP:}} Use \textcolor{myblue}{multiple metrics} - don't rely on just one!
	\end{frame}
	
	\begin{frame}{T74: Test Sample Importance}
		\fontsize{6.5}{7.5}\selectfont
		\textbf{TOPIC: Why Test Sample Comparison is More Important Than Training}
		
		\vspace{0.1cm}
		\textbf{TRAINING SAMPLE:}
		\begin{itemize}
			\item Shows how well model memorizes
			\item Optimistically biased
			\item Not representative of new data
			\item Not a true test of performance
		\end{itemize}
		
		\vspace{0.1cm}
		\textbf{TEST SAMPLE:}
		\begin{itemize}
			\item Shows how well model generalizes
			\item Unbiased performance estimate
			\item What matters in practice
			\item True test of performance
		\end{itemize}
		
		\vspace{0.1cm}
		\textbf{KEY INSIGHT:}
		\begin{itemize}
			\item We are generally more interested in validation or test sample
			\item Models overfit to training sample
			\item Examining accuracy on seen data is not true test
		\end{itemize}
		
		\textcolor{myred}{\textbf{EXAM TRAP:}} Test sample = \textcolor{myblue}{true generalization} - always prioritize!
	\end{frame}
	
	\begin{frame}{T75: Model Interpretation Tradeoff}
		\fontsize{6.5}{7.5}\selectfont
		\textbf{TOPIC: Tradeoff Between Logistic Regression and Neural Networks}
		
		\vspace{0.1cm}
		\textbf{LOGISTIC REGRESSION:}
		\begin{itemize}
			\item \textcolor{myblue}{Pros:} Interpretable, transparent, regulatory friendly
			\item \textcolor{myblue}{Cons:} Limited to linear relationships
		\end{itemize}
		
		\vspace{0.1cm}
		\textbf{NEURAL NETWORKS:}
		\begin{itemize}
			\item \textcolor{myblue}{Pros:} Can capture complex patterns
			\item \textcolor{myblue}{Cons:} Black box, hard to explain
		\end{itemize}
		
		\vspace{0.1cm}
		\textbf{CHOICE DEPENDS ON:}
		\begin{itemize}
			\item Regulatory requirements
			\item Business needs
			\item Stakeholder expectations
			\item Interpretability vs performance
		\end{itemize}
		
		\vspace{0.1cm}
		\textbf{S\&P GLOBAL STUDY:}
		\begin{itemize}
			\item Identified tradeoffs between interpretability and performance
			\item Emphasized choosing models aligned with business objectives
		\end{itemize}
		
		\textcolor{myred}{\textbf{EXAM TRAP:}} Tradeoff = \textcolor{myblue}{interpretability vs performance}!
	\end{frame}
	
	\begin{frame}{T76: LendingClub Key Takeaway}
		\fontsize{6.5}{7.5}\selectfont
		\textbf{TOPIC: Key Takeaway from LendingClub Model Comparison}
		
		\vspace{0.1cm}
		\textbf{KEY INSIGHTS:}
		\begin{itemize}
			\item Different metrics gave different rankings
			\item Logistic regression: Better accuracy and precision
			\item Neural network: Better recall
			\item No clear winner
		\end{itemize}
		
		\vspace{0.1cm}
		\textbf{IMPLICATION:}
		\begin{itemize}
			\item Model choice must align with business objectives
			\item "Contradictory indicators illustrate importance of choosing appropriate evaluation metric"
			\item Regulatory context matters
			\item No universal best model
		\end{itemize}
		
		\vspace{0.1cm}
		\textbf{PRACTICAL ADVICE:}
		\begin{itemize}
			\item Understand what matters most
			\item Choose metrics accordingly
			\item Document rationale
		\end{itemize}
		
		\textcolor{myred}{\textbf{EXAM TRAP:}} No universal best model - \textcolor{myblue}{choice depends on criteria}!
	\end{frame}
	
	\begin{frame}{T77: Business Metric Alignment}
		\fontsize{6.5}{7.5}\selectfont
		\textbf{TOPIC: Why Evaluation Metrics Should Align with Business Objectives}
		
		\vspace{0.1cm}
		\textbf{WHY ALIGNMENT MATTERS:}
		\begin{itemize}
			\item Metrics drive model selection
			\item Model selection drives decisions
			\item Decisions have real-world consequences
			\item Wrong metric = wrong decisions
		\end{itemize}
		
		\vspace{0.1cm}
		\textbf{EXAMPLES:}
		\begin{itemize}
			\item If catching fraud is priority → optimize recall
			\item If minimizing investigations → optimize precision
			\item Must match business needs
		\end{itemize}
		
		\vspace{0.1cm}
		\textbf{ALIGNMENT PROCESS:}
		\begin{enumerate}
			\item Identify business objectives
			\item Determine cost of errors
			\item Choose appropriate metrics
			\item Align threshold with risk appetite
		\end{enumerate}
		
		\textcolor{myred}{\textbf{EXAM TRAP:}} Metrics must \textcolor{myblue}{align with business objectives}!
	\end{frame}
	
	\begin{frame}{T78: Regulatory Considerations}
		\fontsize{6.5}{7.5}\selectfont
		\textbf{TOPIC: How Regulatory Requirements Affect Model Evaluation}
		
		\vspace{0.1cm}
		\textbf{REGULATORY EXPECTATIONS:}
		\begin{itemize}
			\item ECB recommends tail-risk metrics
			\item RMSE alongside MAE
			\item Documentation of metric choice
			\item Justification for model selection
		\end{itemize}
		
		\vspace{0.1cm}
		\textbf{WHY IT MATTERS:}
		\begin{itemize}
			\item Compliance is mandatory
			\item Penalties for non-compliance
			\item Reputational impact
			\item Regulatory scrutiny
		\end{itemize}
		
		\vspace{0.1cm}
		\textbf{KEY REGULATORY REQUIREMENTS:}
		\begin{itemize}
			\item Metrics capturing tail risk
			\item RMSE alongside MAE
			\item Documentation and transparency
			\item Periodic review
		\end{itemize}
		
		\textcolor{myred}{\textbf{EXAM TRAP:}} Regulators \textcolor{myblue}{specify metric requirements}!
	\end{frame}
	
	\begin{frame}{T79: Multiple Model Comparison}
		\fontsize{6.5}{7.5}\selectfont
		\textbf{TOPIC: Framework for Comparing Multiple Models}
		
		\vspace{0.1cm}
		\textbf{COMPARISON FRAMEWORK:}
		\begin{enumerate}
			\item \textcolor{myblue}{Multiple Metrics:}
			\begin{itemize}
				\item Accuracy, Precision, Recall, F1, AUC
				\item Continuous: MSE, RMSE, MAE, MAPE
			\end{itemize}
			
			\item \textcolor{myblue}{Business Context:}
			\begin{itemize}
				\item Costs, priorities, risk appetite
			\end{itemize}
			
			\item \textcolor{myblue}{Regulatory Requirements:}
			\begin{itemize}
				\item Required metrics, documentation
			\end{itemize}
			
			\item \textcolor{myblue}{Interpretability:}
			\begin{itemize}
				\item Stakeholder needs
			\end{itemize}
			
			\item \textcolor{myblue}{Performance:}
			\begin{itemize}
				\item Training vs test gap
			\end{itemize}
		\end{enumerate}
		
		\textcolor{myred}{\textbf{EXAM TRAP:}} Comprehensive comparison needs multiple perspectives!
	\end{frame}
	
	\begin{frame}{T80: Evaluation Framework Summary}
		\fontsize{6.5}{7.5}\selectfont
		\textbf{TOPIC: Recommended Framework for Model Performance Evaluation}
		
		\vspace{0.1cm}
		\textbf{RECOMMENDED FRAMEWORK:}
		\begin{enumerate}
			\item \textcolor{myblue}{Continuous Metrics:}
			\begin{itemize}
				\item MSE, RMSE, MAE, MAPE
			\end{itemize}
			
			\item \textcolor{myblue}{Classification Metrics:}
			\begin{itemize}
				\item Confusion Matrix, Accuracy, Precision, Recall, F1, AUC
			\end{itemize}
			
			\item \textcolor{myblue}{Business Context:}
			\begin{itemize}
				\item Cost matrix, risk appetite
			\end{itemize}
			
			\item \textcolor{myblue}{Documentation:}
			\begin{itemize}
				\item Transparent and defensible
			\end{itemize}
		\end{enumerate}
		
		\vspace{0.1cm}
		\textbf{KEY PRINCIPLE:}
		\begin{itemize}
			\item No single metric tells the whole story
			\item Choose based on problem and stakeholders
			\item Document and justify choices
		\end{itemize}
		
		\textcolor{myred}{\textbf{EXAM TRAP:}} Comprehensive evaluation = \textcolor{myblue}{multiple metrics + business context + documentation}!
	\end{frame}
	
	% ============================================================
	% SECTION 8: ADVANCED TOPICS
	% ============================================================
	
	\begin{frame}{T81: Error Distribution Analysis}
		\fontsize{6.5}{7.5}\selectfont
		\textbf{TOPIC: What Error Distributions Reveal About Model Performance}
		
		\vspace{0.1cm}
		\textbf{KEY INSIGHTS FROM DISTRIBUTIONS:}
		\begin{itemize}
			\item \textcolor{myblue}{Bias:} Systematic over/under prediction
			\item \textcolor{myblue}{Segments:} Areas where model performs poorly
			\item \textcolor{myblue}{Shape:} Many small vs few large errors
		\end{itemize}
		
		\vspace{0.1cm}
		\textbf{VISUALIZATIONS TO USE:}
		\begin{itemize}
			\item Histograms of residuals
			\item Box plots by segment
			\item Q-Q plots for normality
		\end{itemize}
		
		\vspace{0.1cm}
		\textbf{KEY QUESTIONS:}
		\begin{itemize}
			\item "Are errors generally small with a few large ones or consistently moderate?"
			\item "Are errors randomly distributed or systematically biased?"
			\item "Are there segments where errors are consistently larger?"
		\end{itemize}
		
		\textcolor{myred}{\textbf{EXAM TRAP:}} Error distributions reveal \textcolor{myblue}{patterns and biases} that averages hide!
	\end{frame}
	
	\begin{frame}{T82: Systematic Bias Detection}
		\fontsize{6.5}{7.5}\selectfont
		\textbf{TOPIC: What Systematic Bias in Model Errors Indicates}
		
		\vspace{0.1cm}
		\textbf{SYSTEMATIC BIAS EXAMPLES:}
		\begin{itemize}
			\item \textcolor{myblue}{Over-prediction:} Consistently predicting higher values
			\item \textcolor{myblue}{Under-prediction:} Consistently predicting lower values
			\item \textcolor{myblue}{Segment Bias:} Errors in specific subgroups
		\end{itemize}
		
		\vspace{0.1cm}
		\textbf{WHY IT MATTERS:}
		\begin{itemize}
			\item Can lead to systematic business impact
			\item Harder to correct than random errors
			\item May indicate model misspecification
			\item More problematic than random errors even if average is small
		\end{itemize}
		
		\vspace{0.1cm}
		\textbf{EXAMPLE:}
		\begin{itemize}
			\item Model consistently over-predicts recovery rates
			\item Leads to underestimation of losses
			\item Capital shortfall
		\end{itemize}
		
		\textcolor{myred}{\textbf{EXAM TRAP:}} Systematic bias = \textcolor{myblue}{consistent error pattern}!
	\end{frame}
	
	\begin{frame}{T83: Segment Analysis}
		\fontsize{6.5}{7.5}\selectfont
		\textbf{TOPIC: Why Analyze Performance Across Different Segments}
		
		\vspace{0.1cm}
		\textbf{BENEFITS OF SEGMENT ANALYSIS:}
		\begin{itemize}
			\item Identifies model weaknesses
			\item Reveals fairness issues
			\item Highlights data quality problems
			\item Guides model improvements
		\end{itemize}
		
		\vspace{0.1cm}
		\textbf{EXAMPLE SEGMENTS:}
		\begin{itemize}
			\item Geographic regions
			\item Product types
			\item Customer segments
			\item Economic conditions
		\end{itemize}
		
		\vspace{0.1cm}
		\textbf{WHAT TO LOOK FOR:}
		\begin{itemize}
			\item Segments with consistently larger errors
			\item Segments with systematic bias
			\item Segments with different performance patterns
		\end{itemize}
		
		\textcolor{myred}{\textbf{EXAM TRAP:}} Segment analysis reveals \textcolor{myblue}{where models fail}!
	\end{frame}
	
	\begin{frame}{T84: Risk Metrics Regulatory Preference}
		\fontsize{6.5}{7.5}\selectfont
		\textbf{TOPIC: Why Regulators Prefer Metrics Capturing Tail Risk}
		
		\vspace{0.1cm}
		\textbf{TAIL RISK IMPORTANCE:}
		\begin{itemize}
			\item Extreme events cause largest losses
			\item Capital requirements driven by tail risk
			\item Regulatory focus on solvency
			\item Systemic risk concerns
		\end{itemize}
		
		\vspace{0.1cm}
		\textbf{WHY MSE/RMSE:}
		\begin{itemize}
			\item Heavily penalize large errors
			\item Capture tail events better than MAE
			\item Align with regulatory capital frameworks
		\end{itemize}
		
		\vspace{0.1cm}
		\textbf{ECB RECOMMENDATIONS:}
		\begin{itemize}
			\item Metrics capturing tail risk
			\item RMSE alongside MAE
			\item Ensures large, infrequent losses not overlooked
		\end{itemize}
		
		\textcolor{myred}{\textbf{EXAM TRAP:}} Regulators want \textcolor{myblue}{tail-risk metrics} for capital adequacy!
	\end{frame}
	
	\begin{frame}{T85: Cost Matrix Components}
		\fontsize{6.5}{7.5}\selectfont
		\textbf{TOPIC: What Should Be Included in a Cost Matrix}
		
		\vspace{0.1cm}
		\textbf{COST COMPONENTS:}
		\begin{enumerate}
			\item \textcolor{myblue}{Financial:} Losses from missed risks, capital impact
			\item \textcolor{myblue}{Operational:} Investigation costs, resource waste
			\item \textcolor{myblue}{Reputational:} Customer trust, market confidence
			\item \textcolor{myblue}{Regulatory:} Fines, sanctions, scrutiny
		\end{enumerate}
		
		\vspace{0.1cm}
		\textbf{BENEFIT COMPONENTS:}
		\begin{itemize}
			\item \textcolor{myblue}{True Positive:} Losses avoided, risks managed
			\item \textcolor{myblue}{True Negative:} Efficient processing, satisfied customers
		\end{itemize}
		
		\vspace{0.1cm}
		\textbf{KEY INSIGHT:}
		\begin{itemize}
			\item Cost matrix should capture all consequences
			\item Financial and non-financial
			\item Tangible and intangible
		\end{itemize}
		
		\textcolor{myred}{\textbf{EXAM TRAP:}} Cost matrix should capture \textcolor{myblue}{all consequences} - financial and non-financial!
	\end{frame}
	
	\begin{frame}{T86: Threshold Justification}
		\fontsize{6.5}{7.5}\selectfont
		\textbf{TOPIC: How to Justify Threshold Choices to Regulators}
		
		\vspace{0.1cm}
		\textbf{JUSTIFICATION SHOULD INCLUDE:}
		\begin{enumerate}
			\item Analysis of tradeoffs
			\item Cost-benefit calculations
			\item Alignment with risk appetite
			\item Stakeholder consultation
			\item Periodic review process
		\end{enumerate}
		
		\vspace{0.1cm}
		\textbf{EXAMPLE STATEMENT:}
		\begin{itemize}
			\item "Threshold chosen to achieve 80\% recall given severe consequences of missing fraud events"
			\item "Accepts X\% false positives costing Y in operational overhead"
		\end{itemize}
		
		\vspace{0.1cm}
		\textbf{WHY IT MATTERS:}
		\begin{itemize}
			\item Regulatory compliance
			\item Audit trail
			\item Defensible decisions
		\end{itemize}
		
		\textcolor{myred}{\textbf{EXAM TRAP:}} Thresholds must be \textcolor{myblue}{justified with documented analysis}!
	\end{frame}
	
	\begin{frame}{T87: Stakeholder Communication}
		\fontsize{6.5}{7.5}\selectfont
		\textbf{TOPIC: How to Communicate Model Performance to Different Stakeholders}
		
		\vspace{0.1cm}
		\textbf{AUDIENCE CONSIDERATIONS:}
		\begin{itemize}
			\item \textcolor{myblue}{Board/Management:} MAPE, Accuracy (percentage-based)
			\item \textcolor{myblue}{Regulators:} RMSE with MAE, tail-risk metrics
			\item \textcolor{myblue}{Model Developers:} All metrics, error distributions
			\item \textcolor{myblue}{Audit Committee:} MAPE for standardized comparison
		\end{itemize}
		
		\vspace{0.1cm}
		\textbf{COMMUNICATION BEST PRACTICES:}
		\begin{itemize}
			\item Use visualizations (ROC curves, error distributions)
			\item Explain what metrics mean in business terms
			\item Highlight tradeoffs and limitations
		\end{itemize}
		
		\textcolor{myred}{\textbf{EXAM TRAP:}} Different stakeholders need \textcolor{myblue}{different metrics and explanations}!
	\end{frame}
	
	\begin{frame}{T88: Performance Monitoring Frequency}
		\fontsize{6.5}{7.5}\selectfont
		\textbf{TOPIC: How Often Should Model Performance Be Monitored?}
		
		\vspace{0.1cm}
		\textbf{MONITORING FREQUENCY:}
		\begin{itemize}
			\item \textcolor{myblue}{Regular Intervals:} Quarterly, semi-annually, annually
			\item \textcolor{myblue}{Event-Driven:} Significant environmental changes
			\item \textcolor{myblue}{Trigger-Based:} Performance threshold breaches
		\end{itemize}
		
		\vspace{0.1cm}
		\textbf{WHAT TO MONITOR:}
		\begin{itemize}
			\item Performance metrics over time
			\item Error distributions
			\item Segment performance
			\item Threshold appropriateness
		\end{itemize}
		
		\vspace{0.1cm}
		\textbf{WHY CONTINUOUS MONITORING:}
		\begin{itemize}
			\item Detect performance decay
			\item Identify concept drift
			\item Ensure ongoing compliance
		\end{itemize}
		
		\textcolor{myred}{\textbf{EXAM TRAP:}} Model performance must be \textcolor{myblue}{monitored continuously} - not just during development!
	\end{frame}
	
	\begin{frame}{T89: Model Validation Process}
		\fontsize{6.5}{7.5}\selectfont
		\textbf{TOPIC: Role of Validation in Model Performance Evaluation}
		
		\vspace{0.1cm}
		\textbf{VALIDATION PURPOSES:}
		\begin{enumerate}
			\item Independent performance assessment
			\item Hyperparameter tuning
			\item Overfitting detection
			\item Model comparison
		\end{enumerate}
		
		\vspace{0.1cm}
		\textbf{VALIDATION BEST PRACTICES:}
		\begin{itemize}
			\item Use separate validation data
			\item Use cross-validation
			\item Document validation process
			\item Regulatory compliance
		\end{itemize}
		
		\vspace{0.1cm}
		\textbf{KEY INSIGHT:}
		\begin{itemize}
			\item Validation provides independent assessment
			\item Different from training
			\item Essential for model governance
		\end{itemize}
		
		\textcolor{myred}{\textbf{EXAM TRAP:}} Validation provides \textcolor{myblue}{independent assessment} of model performance!
	\end{frame}
	
	\begin{frame}{T90: Performance Metrics Documentation}
		\fontsize{6.5}{7.5}\selectfont
		\textbf{TOPIC: Why Performance Metric Choices Should Be Documented}
		
		\vspace{0.1cm}
		\textbf{DOCUMENTATION SHOULD INCLUDE:}
		\begin{enumerate}
			\item Choice of metrics
			\item Justification for choices
			\item Business context
			\item Regulatory requirements
			\item Performance results
			\item Model selection rationales
		\end{enumerate}
		
		\vspace{0.1cm}
		\textbf{WHY IT MATTERS:}
		\begin{itemize}
			\item Regulatory compliance
			\item Audit trail
			\item Knowledge transfer
			\item Defensible decisions
		\end{itemize}
		
		\textcolor{myred}{\textbf{EXAM TRAP:}} Metric choices must be \textcolor{myblue}{documented and justified}!
	\end{frame}
	
	\begin{frame}{T91: Performance Decay Detection}
		\fontsize{6.5}{7.5}\selectfont
		\textbf{TOPIC: How to Detect Model Performance Decay}
		
		\vspace{0.1cm}
		\textbf{DETECTION METHODS:}
		\begin{itemize}
			\item Track metrics over time
			\item Compare to baseline performance
			\item Monitor error distributions
			\item Segment performance tracking
		\end{itemize}
		
		\vspace{0.1cm}
		\textbf{COMMON CAUSES:}
		\begin{itemize}
			\item Data drift (concept drift)
			\item Economic changes
			\item Regulatory changes
			\item Model degradation
		\end{itemize}
		
		\vspace{0.1cm}
		\textbf{WHY IT MATTERS:}
		\begin{itemize}
			\item Early detection allows remediation
			\item Prevents silent model failure
			\item Maintains performance standards
		\end{itemize}
		
		\textcolor{myred}{\textbf{EXAM TRAP:}} Monitor metrics over time to detect \textcolor{myblue}{performance decay}!
	\end{frame}
	
	\begin{frame}{T92: Concept Drift}
		\fontsize{6.5}{7.5}\selectfont
		\textbf{TOPIC: What is Concept Drift and How Does It Affect Performance?}
		
		\vspace{0.1cm}
		\textbf{DEFINITION:}
		\begin{itemize}
			\item Changes in underlying relationship between inputs and outputs over time
			\item Model trained on historical data may become less accurate
		\end{itemize}
		
		\vspace{0.1cm}
		\textbf{EXAMPLES:}
		\begin{itemize}
			\item Changing fraud patterns
			\item Economic cycles affecting default relationships
			\item Customer behavior changes
			\item Regulatory changes
		\end{itemize}
		
		\vspace{0.1cm}
		\textbf{DETECTION AND RESPONSE:}
		\begin{itemize}
			\item Monitor performance metrics
			\item Regularly retrain models
			\item Update model parameters
			\item Review thresholds
		\end{itemize}
		
		\textcolor{myred}{\textbf{EXAM TRAP:}} Concept drift = \textcolor{myblue}{changing relationships} over time - requires monitoring!
	\end{frame}
	
	\begin{frame}{T93: Model Retraining Triggers}
		\fontsize{6.5}{7.5}\selectfont
		\textbf{TOPIC: What Should Trigger Model Retraining?}
		
		\vspace{0.1cm}
		\textbf{RETRAINING TRIGGERS:}
		\begin{enumerate}
			\item \textcolor{myblue}{Performance:} Metrics below acceptable levels
			\item \textcolor{myblue}{Data:} Significant changes in data distribution
			\item \textcolor{myblue}{Schedule:} Regular intervals (monthly, quarterly)
			\item \textcolor{myblue}{Events:} Major economic or regulatory changes
		\end{enumerate}
		
		\vspace{0.1cm}
		\textbf{BEST PRACTICE:}
		\begin{itemize}
			\item Document retraining triggers
			\item Regular performance review
			\item Proactive monitoring
			\item Automated alerts
		\end{itemize}
		
		\vspace{0.1cm}
		\textbf{KEY INSIGHT:}
		\begin{itemize}
			\item Models are not permanent
			\item Need periodic refresh
			\item Align with business changes
		\end{itemize}
		
		\textcolor{myred}{\textbf{EXAM TRAP:}} Retrain when \textcolor{myblue}{performance decays, data changes, or on schedule}!
	\end{frame}
	
	\begin{frame}{T94: Performance Reporting}
		\fontsize{6.5}{7.5}\selectfont
		\textbf{TOPIC: What Should Be Included in Model Performance Reports}
		
		\vspace{0.1cm}
		\textbf{REPORT COMPONENTS:}
		\begin{enumerate}
			\item \textcolor{myblue}{Current:} Latest performance metrics
			\item \textcolor{myblue}{Trends:} Performance over time
			\item \textcolor{myblue}{Baseline:} Comparison to initial/expected performance
			\item \textcolor{myblue}{Issues:} Any identified problems or concerns
			\item \textcolor{myblue}{Actions:} Planned or taken remediation
		\end{enumerate}
		
		\vspace{0.1cm}
		\textbf{AUDIENCE:}
		\begin{itemize}
			\item Model owners
			\item Management
			\item Regulators
			\item Audit committees
		\end{itemize}
		
		\vspace{0.1cm}
		\textbf{KEY INSIGHT:}
		\begin{itemize}
			\item Reports should include current metrics, trends, and issues
			\item Regular and transparent
		\end{itemize}
		
		\textcolor{myred}{\textbf{EXAM TRAP:}} Reports should include \textcolor{myblue}{current metrics, trends, and issues}!
	\end{frame}
	
	\begin{frame}{T95: Fairness and Bias}
		\fontsize{6.5}{7.5}\selectfont
		\textbf{TOPIC: How Model Performance Evaluation Helps Identify Fairness Issues}
		
		\vspace{0.1cm}
		\textbf{FAIRNESS EVALUATION:}
		\begin{itemize}
			\item Compare performance across segments
			\item Check for systematic bias
			\item Ensure equal treatment
			\item Regulatory compliance (fair lending)
		\end{itemize}
		
		\vspace{0.1cm}
		\textbf{WHAT TO EXAMINE:}
		\begin{itemize}
			\item Accuracy differences
			\item False positive/negative rates
			\item Disparate impact
			\item Equal opportunity
		\end{itemize}
		
		\vspace{0.1cm}
		\textbf{WHY IT MATTERS:}
		\begin{itemize}
			\item Legal and regulatory requirements
			\item Ethical considerations
			\item Reputational risk
		\end{itemize}
		
		\textcolor{myred}{\textbf{EXAM TRAP:}} Segment analysis reveals \textcolor{myblue}{fairness issues}!
	\end{frame}
	
	\begin{frame}{T96: Explainable AI}
		\fontsize{6.5}{7.5}\selectfont
		\textbf{TOPIC: How Model Performance Evaluation Relates to Explainable AI}
		
		\vspace{0.1cm}
		\textbf{CONNECTION TO XAI:}
		\begin{itemize}
			\item Poor performance in segments → need explanations
			\item Complex models (DNNs) → need XAI techniques
			\item Regulatory requirements for explanation
			\item Tools: SHAP, LIME for interpretation
		\end{itemize}
		
		\vspace{0.1cm}
		\textbf{PRACTICAL APPLICATION:}
		\begin{itemize}
			\item Use explanations to understand poor performance
			\item Identify feature importance
			\item Debug model issues
			\item Build trust in predictions
		\end{itemize}
		
		\textcolor{myred}{\textbf{EXAM TRAP:}} Performance evaluation guides \textcolor{myblue}{where explanations are needed}!
	\end{frame}
	
	\begin{frame}{T97: SHAP Values}
		\fontsize{6.5}{7.5}\selectfont
		\textbf{TOPIC: How SHAP Values Help with Model Performance Evaluation}
		
		\vspace{0.1cm}
		\textbf{SHAP IN EVALUATION:}
		\begin{itemize}
			\item Identify important features
			\item Understand model behavior
			\item Debug performance issues
			\item Explain to stakeholders
		\end{itemize}
		
		\vspace{0.1cm}
		\textbf{APPLICATION:}
		\begin{itemize}
			\item Feature importance analysis
			\item Understanding model decisions
			\item Regulatory explanations
			\item Model improvement insights
		\end{itemize}
		
		\vspace{0.1cm}
		\textbf{KEY INSIGHT:}
		\begin{itemize}
			\item SHAP helps understand what drives model performance
			\item Identifies which features are driving good and bad performance
		\end{itemize}
		
		\textcolor{myred}{\textbf{EXAM TRAP:}} SHAP helps \textcolor{myblue}{understand what drives model performance}!
	\end{frame}
	
	\begin{frame}{T98: LIME for Local Explanations}
		\fontsize{6.5}{7.5}\selectfont
		\textbf{TOPIC: When LIME is Useful for Model Evaluation}
		
		\vspace{0.1cm}
		\textbf{LIME APPLICATIONS:}
		\begin{itemize}
			\item Explain specific predictions
			\item Understand why model made certain decisions
			\item Identify problematic cases
			\item Build trust with stakeholders
		\end{itemize}
		
		\vspace{0.1cm}
		\textbf{EVALUATION INSIGHTS:}
		\begin{itemize}
			\item Why certain cases were misclassified
			\item What features drove errors
			\item Edge case behavior
			\item Model limitations
		\end{itemize}
		
		\vspace{0.1cm}
		\textbf{KEY INSIGHT:}
		\begin{itemize}
			\item LIME explains individual predictions
			\item Useful for error analysis
			\item Complements global metrics
		\end{itemize}
		
		\textcolor{myred}{\textbf{EXAM TRAP:}} LIME explains \textcolor{myblue}{individual predictions} - useful for error analysis!
	\end{frame}
	
	\begin{frame}{T99: Model Governance}
		\fontsize{6.5}{7.5}\selectfont
		\textbf{TOPIC: Role of Model Governance in Performance Evaluation}
		
		\vspace{0.1cm}
		\textbf{GOVERNANCE COMPONENTS:}
		\begin{enumerate}
			\item \textcolor{myblue}{Standardization:} Consistent metrics across models
			\item \textcolor{myblue}{Documentation:} Clear records of evaluation
			\item \textcolor{myblue}{Monitoring:} Regular performance review
			\item \textcolor{myblue}{Oversight:} Independent validation
		\end{enumerate}
		
		\vspace{0.1cm}
		\textbf{WHY IT MATTERS:}
		\begin{itemize}
			\item Regulatory compliance
			\item Risk management
			\item Consistent decision-making
			\item Auditability
		\end{itemize}
		
		\vspace{0.1cm}
		\textbf{KEY INSIGHT:}
		\begin{itemize}
			\item Governance ensures consistent evaluation and monitoring
			\item Essential for model risk management
		\end{itemize}
		
		\textcolor{myred}{\textbf{EXAM TRAP:}} Governance ensures \textcolor{myblue}{consistent evaluation and monitoring}!
	\end{frame}
	
	\begin{frame}{T100: Chapter Summary}
		\fontsize{6.5}{7.5}\selectfont
		\textbf{TOPIC: Key Takeaways from Model Performance Evaluation}
		
		\vspace{0.1cm}
		\textbf{KEY TAKEAWAYS:}
		\begin{enumerate}
			\item \textcolor{myblue}{Multiple Metrics:} Use continuous (MSE, RMSE, MAE, MAPE) and classification (Confusion Matrix, Accuracy, Precision, Recall, F1, AUC)
			\item \textcolor{myblue}{Business Context:} Align with business objectives and risk appetite
			\item \textcolor{myblue}{Continuous Monitoring:} Track performance over time
			\item \textcolor{myblue}{Documentation:} Record and justify all choices
		\end{enumerate}
		
		\vspace{0.1cm}
		\textbf{CRITICAL INSIGHTS:}
		\begin{itemize}
			\item No single metric tells the whole story
			\item Different metrics serve different purposes
			\item Performance evaluation is a continuous process
			\item Regulatory requirements must be met
		\end{itemize}
		
		\vspace{0.1cm}
		\textbf{FINAL ADVICE:}
		\begin{itemize}
			\item Understand metrics, tradeoffs, and when to use each
			\item Choose based on problem and stakeholders
			\item Document and justify choices
		\end{itemize}
		
		\textcolor{myred}{\textbf{EXAM TRAP:}} Comprehensive evaluation = \textcolor{myblue}{multiple metrics + business context + monitoring + documentation}!
	\end{frame}
	
	% ============================================================
	% SUMMARY SLIDE
	% ============================================================
	
	\begin{frame}{Summary: Complete Review}
		\fontsize{6.5}{7.5}\selectfont
		\textbf{\textcolor{myblue}{Continuous Output Metrics:}}
		\begin{itemize}
			\item \textcolor{myblue}{MSE:} Average squared error; penalizes large errors; captures tail risk
			\item \textcolor{myblue}{RMSE:} Square root of MSE; same units as target; easier to interpret
			\item \textcolor{myblue}{MAE:} Average absolute error; robust to outliers
			\item \textcolor{myblue}{MAPE:} Percentage error; board-friendly; compares across scales
		\end{itemize}
		
		\vspace{0.2cm}
		\textbf{\textcolor{myblue}{Classification Metrics:}}
		\begin{itemize}
			\item \textcolor{myblue}{Confusion Matrix:} TP, TN, FP, FN - foundation of all metrics
			\item \textcolor{myblue}{Accuracy:} Total correct / total - misleading for imbalanced data
			\item \textcolor{myblue}{Precision:} TP / (TP + FP) - "of predicted positives, how many were right?"
			\item \textcolor{myblue}{Recall:} TP / (TP + FN) - "of actual positives, how many did we catch?"
			\item \textcolor{myblue}{F1:} Harmonic mean of precision and recall
			\item \textcolor{myblue}{AUC:} Area under ROC curve - threshold-independent ranking ability
		\end{itemize}
		
		\vspace{0.2cm}
		\textbf{\textcolor{myblue}{Threshold Selection:}}
		\begin{itemize}
			\item \textcolor{myblue}{Cost Matrix:} Economic consequences of each decision
			\item \textcolor{myblue}{Tradeoff:} Balance false positives vs false negatives
			\item \textcolor{myblue}{Dynamic:} Adjust based on changing conditions
			\item \textcolor{myblue}{Documentation:} Transparent and defensible
		\end{itemize}
		
		\vspace{0.2cm}
		\textcolor{myred}{\textbf{Exam Success:}} Understand the metrics, their tradeoffs, and when to use each!
	\end{frame}
	
\end{document}